% Complete reviewed Unit 7; editorial master: chapters/ch07/answer-key.md
% Markdown SHA-256: 2e70c0efc0dc5eb9c83a8a2a4fa56c4ec1ee7778cf1d5bb65cb9bde82448d0b2
% One-use initial assembly followed by complete content and page review.
\newcommand{\runningtitle}{Unit 7. Answer key}
% Layout inherited from the reviewed previous course; set a running title per unit.
% This is a small hand-maintained TeX layout, not a conversion pipeline.
\RequirePackage{fix-cm}
\documentclass[11pt,twoside]{article}
\usepackage{fontspec}
\setmainfont{Cambria}[
  Renderer=HarfBuzz,
  BoldFont=cambriab.ttf,
  ItalicFont=cambriai.ttf,
  BoldItalicFont=cambriaz.ttf
]
\setmonofont{lmmono10-regular.otf}[BoldFont=lmmonolt10-bold.otf]
\usepackage[paperwidth=165mm,paperheight=241mm,inner=21mm,outer=19mm,
  top=22mm,bottom=20mm,headheight=14pt,headsep=7mm,footskip=10mm]{geometry}
\usepackage{xcolor,array,booktabs,longtable,tabularx}
\usepackage{enumitem,fancyhdr}
\usepackage[unicode,hidelinks]{hyperref}
\definecolor{ink}{HTML}{000000}
\definecolor{muted}{HTML}{4C5558}
\setlength{\parindent}{0pt}
\setlength{\parskip}{3.5pt}
\AtBeginDocument{\fontsize{11}{14.3}\selectfont}
\setlength{\tabcolsep}{4pt}
\renewcommand{\arraystretch}{1.10}
\setlist[itemize]{leftmargin=1.3em,label=\textbullet,itemsep=2pt,topsep=3pt,parsep=0pt}
\setlist[enumerate]{leftmargin=1.8em,itemsep=3pt,topsep=4pt,parsep=0pt}
\widowpenalty=10000
\clubpenalty=10000
\emergencystretch=1em
\raggedbottom
\pagestyle{fancy}
\fancyhf{}
\providecommand{\runningtitle}{Old English}
\fancyhead[LE,RO]{\small\itshape \runningtitle}
\fancyhead[LO,RE]{\small Old English}
\fancyfoot[LE,RO]{\small\thepage}
\renewcommand{\headrulewidth}{0.3pt}
\newcommand{\booktitle}[2]{
  {\fontsize{9}{11}\selectfont\color{muted}#1\par}
  \vspace{3pt}
  {\fontsize{20}{24}\selectfont\color{ink}#2\par}
  \vspace{6pt}
}
\newcommand{\source}[1]{\par{\footnotesize\color{muted}#1\par}}
\newcommand{\readingtext}[1]{
  \par\addvspace{5pt}
  \begingroup\leftskip=1.5em
    \fontsize{12}{16}\selectfont #1\par
  \endgroup\addvspace{5pt}
}
\makeatletter
\renewcommand{\section}{\@startsection{section}{1}{\z@}
  {-2.4ex \@plus -0.6ex \@minus -0.2ex}{0.8ex \@plus 0.2ex}
  {\normalfont\fontsize{12}{15}\selectfont\bfseries\color{ink}}}
\renewcommand{\subsection}{\@startsection{subsection}{2}{\z@}
  {-1.8ex \@plus -0.5ex \@minus -0.2ex}{0.5ex \@plus 0.2ex}
  {\normalfont\fontsize{11}{14}\selectfont\bfseries\color{ink}}}
\makeatother

\hyphenpenalty=10000
\exhyphenpenalty=10000
\emergencystretch=3em
\newcommand{\itempara}[2]{\par\addvspace{3pt}\noindent\hangindent=1.8em\hangafter=1\makebox[1.8em][l]{#1.}#2\par}
\newcommand{\reserve}[1]{\par\penalty-100\begingroup\dimen0=\pagegoal\advance\dimen0 by-\pagetotal\ifdim\dimen0<#1\newpage\fi\endgroup}
\begin{document}
\addfontfeatures{Ligatures=TeXOff}
\raggedright
\booktitle{UNIT 7}{Answer key}

% mdparagraph:P001:begin
The models below give the selected forms of this course. A different word order or an attested form can be good Old English; assess it separately from compliance with an expressly requested verb, construction, or paradigm. Preserve the meaning, participants, number, tense, and force of the task. A spelling shared by several forms does not by itself decide its analysis.
% mdparagraph:P001:end

% mdparagraph:P002:begin
In the source reading, the printed help remains available. The three forms \textbf{swimmendu, flēogendu, libbendu} are for recognition in that passage. The supplied clauses and exceptional forms are not tests of unaided recall.
% mdparagraph:P002:end

\reserve{7\baselineskip}
% mdparagraph:P003:begin
\section*{A. Subjunctive forms}
% mdparagraph:P003:end

% mdparagraph:P004:begin
In A01–A05, every singular entry serves the first, second, and third persons singular; every plural entry serves the first, second, and third persons plural. Subjects identify person. A dual subject likewise takes a plural verb. The four displayed forms for each verb therefore give the complete present and preterite subjunctive systems requested. [1]
% mdparagraph:P004:end

\par\begin{minipage}{\linewidth}\raggedright
% item: NEW-U07-A01
% key: KEY-NEW-U07-A01
% mdparagraph:P005:begin
\textbf{A01. hīeran and sęllan.}
% mdparagraph:P005:end

% mdtable:T01:begin
% course-table-style: 2026-10
% course-panel:T01-P01:begin
\par\begingroup
\fontsize{10.5}{13.2}\selectfont
\setlength{\LTleft}{0pt}\setlength{\LTright}{\fill}
\noindent
\begin{tabular}{@{}>{\raggedright\arraybackslash}p{20.873mm}>{\raggedright\arraybackslash}p{21.064mm}>{\raggedright\arraybackslash}p{24.876mm}>{\raggedright\arraybackslash}p{21.064mm}>{\raggedright\arraybackslash}p{24.876mm}@{}}
\toprule
\textbf{Person and number} & \textbf{hīeran · Present} & \textbf{hīeran · Preterite} & \textbf{sęllan · Present} & \textbf{sęllan · Preterite} \\
\midrule
1–3 sg. & hīere & hīerde & sęlle & sealde \\
1–3 pl. & hīeren & hīerden & sęllen & sealden \\
\bottomrule
\end{tabular}
\par\endgroup
% course-panel:T01-P01:end
% mdtable:T01:end
% mdparagraph:P006:begin
\textbf{Target:} distinguish the ordinary weak-I dental preterite from the lexical preterite of \textbf{sęllan}. Retain \textbf{hīer-} in both tenses; use \textbf{seald-}, not an invented dental extension of \textbf{sęll-}, in the preterite. The plural subjunctive ends in \textbf{-en}; \textbf{hīerdon, sealdon} would be indicative.
% mdparagraph:P006:end

\end{minipage}\par\addvspace{4pt}

\par\begin{minipage}{\linewidth}\raggedright
% item: NEW-U07-A02
% key: KEY-NEW-U07-A02
% mdparagraph:P007:begin
\textbf{A02. lufian and libban.}
% mdparagraph:P007:end

% mdtable:T02:begin
% course-table-style: 2026-10
% course-panel:T02-P01:begin
\par\begingroup
\fontsize{10.5}{13.2}\selectfont
\setlength{\LTleft}{0pt}\setlength{\LTright}{\fill}
\noindent
\begin{tabular}{@{}>{\raggedright\arraybackslash}p{20.873mm}>{\raggedright\arraybackslash}p{21.064mm}>{\raggedright\arraybackslash}p{24.876mm}>{\raggedright\arraybackslash}p{21.064mm}>{\raggedright\arraybackslash}p{24.876mm}@{}}
\toprule
\textbf{Person and number} & \textbf{lufian · Present} & \textbf{lufian · Preterite} & \textbf{libban · Present} & \textbf{libban · Preterite} \\
\midrule
1–3 sg. & lufie & lufode & libbe & lifde \\
1–3 pl. & lufien & lufoden & libben & lifden \\
\bottomrule
\end{tabular}
\par\endgroup
% course-panel:T02-P01:end
% mdtable:T02:end
% mdparagraph:P008:begin
\textbf{Target:} weak-II \textbf{-ie/\allowbreak{}-ien} and the selected weak-III stems. The \textbf{i} of \textbf{lufi-} remains before the present subjunctive endings. For \textbf{libban}, distinguish \textbf{libb-} here from present indicative \textbf{leof-} in \textbf{leofað}, and from preterite \textbf{lifd-}. Do not make the subjunctive by retaining the indicative third-singular ending \textbf{-að}.
% mdparagraph:P008:end

\end{minipage}\par\addvspace{4pt}

\par\begin{minipage}{\linewidth}\raggedright
% item: NEW-U07-A03
% key: KEY-NEW-U07-A03
% mdparagraph:P009:begin
\textbf{A03. bindan and weaxan.}
% mdparagraph:P009:end

% mdtable:T03:begin
% course-table-style: 2026-10
% course-panel:T03-P01:begin
\par\begingroup
\fontsize{10.5}{13.2}\selectfont
\setlength{\LTleft}{0pt}\setlength{\LTright}{\fill}
\noindent
\begin{tabular}{@{}>{\raggedright\arraybackslash}p{20.873mm}>{\raggedright\arraybackslash}p{21.064mm}>{\raggedright\arraybackslash}p{24.876mm}>{\raggedright\arraybackslash}p{21.064mm}>{\raggedright\arraybackslash}p{24.876mm}@{}}
\toprule
\textbf{Person and number} & \textbf{bindan · Present} & \textbf{bindan · Preterite} & \textbf{weaxan · Present} & \textbf{weaxan · Preterite} \\
\midrule
1–3 sg. & binde & bunde & weaxe & wēoxe \\
1–3 pl. & binden & bunden & weaxen & wēoxen \\
\bottomrule
\end{tabular}
\par\endgroup
% course-panel:T03-P01:end
% mdtable:T03:end
% mdparagraph:P010:begin
\textbf{Target:} use the strong preterite plural stem. Compare indicative \textbf{bundon, wēoxon}. \textbf{Band} is a singular indicative form and does not supply the stem of \textbf{bunde/\allowbreak{}bunden}. The present subjunctive \textbf{weaxe} also keeps the present stem, rather than the altered vowel of indicative \textbf{wiext}.
% mdparagraph:P010:end

\end{minipage}\par\addvspace{4pt}

\par\begin{minipage}{\linewidth}\raggedright
% item: NEW-U07-A04
% key: KEY-NEW-U07-A04
% mdparagraph:P011:begin
\textbf{A04. sēon and biddan.}
% mdparagraph:P011:end

% mdtable:T04:begin
% course-table-style: 2026-10
% course-panel:T04-P01:begin
\par\begingroup
\fontsize{10.5}{13.2}\selectfont
\setlength{\LTleft}{0pt}\setlength{\LTright}{\fill}
\noindent
\begin{tabular}{@{}>{\raggedright\arraybackslash}p{21.528mm}>{\raggedright\arraybackslash}p{19.955mm}>{\raggedright\arraybackslash}p{23.887mm}>{\raggedright\arraybackslash}p{21.725mm}>{\raggedright\arraybackslash}p{25.657mm}@{}}
\toprule
\textbf{Person and number} & \textbf{sēon · Present} & \textbf{sēon · Preterite} & \textbf{biddan · Present} & \textbf{biddan · Preterite} \\
\midrule
1–3 sg. & sēo & sāwe & bidde & bǣde \\
1–3 pl. & sēon & sāwen & bidden & bǣden \\
\bottomrule
\end{tabular}
\par\endgroup
% course-panel:T04-P01:end
% mdtable:T04:end
% mdparagraph:P012:begin
\textbf{Target:} preserve the contracted present of \textbf{sēon} and the doubled consonant of present \textbf{bidd-}. The preterites use \textbf{sāw-} and \textbf{bǣd-}, corresponding to the learned plural indicative stems. \textbf{Sēon} can also spell the infinitive, but that does not remove its place in the finite subjunctive paradigm; a clause determines which use is present.
% mdparagraph:P012:end

\end{minipage}\par\addvspace{4pt}

\par\begin{minipage}{\linewidth}\raggedright
% item: NEW-U07-A05
% key: KEY-NEW-U07-A05
% mdparagraph:P013:begin
\textbf{A05. witan and willan.}
% mdparagraph:P013:end

% mdtable:T05:begin
% course-table-style: 2026-10
% course-panel:T05-P01:begin
\par\begingroup
\fontsize{10.5}{13.2}\selectfont
\setlength{\LTleft}{0pt}\setlength{\LTright}{\fill}
\noindent
\begin{tabular}{@{}>{\raggedright\arraybackslash}p{21.195mm}>{\raggedright\arraybackslash}p{20.518mm}>{\raggedright\arraybackslash}p{24.389mm}>{\raggedright\arraybackslash}p{21.389mm}>{\raggedright\arraybackslash}p{25.26mm}@{}}
\toprule
\textbf{Person and number} & \textbf{witan · Present} & \textbf{witan · Preterite} & \textbf{willan · Present} & \textbf{willan · Preterite} \\
\midrule
1–3 sg. & wite & wiste & wille & wolde \\
1–3 pl. & witen & wisten & willen & wolden \\
\bottomrule
\end{tabular}
\par\endgroup
% course-panel:T05-P01:end
% mdtable:T05:end
% mdparagraph:P014:begin
\textbf{Target:} retain each irregular system. The present indicative \textbf{wāt} does not supply the subjunctive \textbf{wite}. \textbf{Wolde} can also be a preterite indicative singular; the spelling alone neither proves subjunctive mood nor requires the English word “would.”
% mdparagraph:P014:end

\end{minipage}\par\addvspace{4pt}

\par\begin{minipage}{\linewidth}\raggedright
% item: NEW-U07-A06
% key: KEY-NEW-U07-A06
% mdparagraph:P015:begin
\textbf{A06. wesan/\allowbreak{}bēon.}
% mdparagraph:P015:end

% mdtable:T06:begin
% course-table-style: 2026-10
% course-panel:T06-P01:begin
\par\begingroup
\fontsize{10.5}{13.2}\selectfont
\setlength{\LTleft}{0pt}\setlength{\LTright}{\fill}
\noindent
\begin{tabular}{@{}>{\raggedright\arraybackslash}p{25.825mm}>{\raggedright\arraybackslash}p{30.542mm}>{\raggedright\arraybackslash}p{30.542mm}>{\raggedright\arraybackslash}p{28.655mm}@{}}
\toprule
\textbf{Person and number} & \textbf{Present, sīe series} & \textbf{Present, bēo series} & \textbf{Shared preterite} \\
\midrule
1–3 sg. & sīe & bēo & wǣre \\
1–3 pl. & sīen & bēon & wǣren \\
\bottomrule
\end{tabular}
\par\endgroup
% course-panel:T06-P01:end
% mdtable:T06:end
% mdparagraph:P016:begin
\textbf{Target:} both selected present alternatives and the one shared preterite. Number is distinguished; person is not distinguished by the ending. Do not invent a second preterite for the \textbf{bēo} series. \textbf{Bēo} also occurs as a singular imperative, and \textbf{bēon} as an infinitive; analyse their use in a clause, not from the spelling alone.
% mdparagraph:P016:end

\end{minipage}\par\addvspace{4pt}

\par\begin{minipage}{\linewidth}\raggedright
% item: NEW-U07-A07
% key: KEY-NEW-U07-A07
% mdparagraph:P017:begin
\textbf{A07. Five new verbal systems.}
% mdparagraph:P017:end

% mdtable:T07:begin
% course-table-style: 2026-10
% course-panel:T07-P01:begin
\par\begingroup
\fontsize{10.5}{13.2}\selectfont
\setlength{\LTleft}{0pt}\setlength{\LTright}{\fill}
\noindent
\begin{tabular}{@{}>{\raggedright\arraybackslash}p{25.0mm}>{\raggedright\arraybackslash}p{27.5mm}>{\raggedright\arraybackslash}p{31.0mm}>{\raggedright\arraybackslash}p{32.064mm}@{}}
\toprule
\textbf{Verb} & \textbf{Present indicative, 3 sg.} & \textbf{Preterite indicative, 3 pl.} & \textbf{Present subjunctive, 1–3 pl.} \\
\midrule
gīeman & gīemð & gīemdon & gīemen \\
rīcsian & rīcsað & rīcsodon & rīcsien \\
gemanigfieldan & gemanigfielt & gemanigfieldon & gemanigfielden \\
tōtwǣman & tōtwǣmð & tōtwǣmdon & tōtwǣmen \\
wyrcean & wyrcð & worhton & wyrcen \\
\bottomrule
\end{tabular}
\par\endgroup
% course-panel:T07-P01:end
% mdtable:T07:end
% mdparagraph:P018:begin
\textbf{Target:} distinguish indicative \textbf{-on} from subjunctive \textbf{-en}, while retaining the lexical stem changes. \textbf{Rīcsien} has the weak-II present \textbf{i}. In \textbf{gemanigfielt}, the stem-final dental and the third-singular ending have the selected contracted outcome; \textbf{gemanigfieldon} retains \textbf{d} before the plural ending. \textbf{Worhton} belongs to the irregular dental preterite of \textbf{wyrcean}, not to a new strong-verb pattern.
% mdparagraph:P018:end

\end{minipage}\par\addvspace{4pt}

\par\begin{minipage}{\linewidth}\raggedright
% item: NEW-U07-A08
% key: KEY-NEW-U07-A08
% mdparagraph:P019:begin
\textbf{A08. Indicative and subjunctive endings.}
% mdparagraph:P019:end

% mdtable:T08:begin
\par\begingroup
\fontsize{10.5}{13.2}\selectfont
\setlength{\LTleft}{0pt}\setlength{\LTright}{\fill}
\noindent
\begin{tabular}{@{}>{\raggedright\arraybackslash}p{39.82mm}>{\raggedright\arraybackslash}p{81.37mm}@{}}
\toprule
\textbf{Given expression} & \textbf{Finite analysis} \\
\midrule
hīe bundon & bindan, preterite indicative, third plural \\
hīe bunden & bindan, preterite subjunctive, third plural \\
hīe lufodon & lufian, preterite indicative, third plural \\
hīe lufoden & lufian, preterite subjunctive, third plural \\
\bottomrule
\end{tabular}
\par\endgroup

% mdtable:T08:end
% mdparagraph:P020:begin
\textbf{Target:} the contrast \textbf{-on : -en}. The verb ending distinguishes mood; \textbf{hīe} supplies third-person reference. \textbf{Bunden} is a finite subjunctive in the expression as presented, although the same spelling also belongs to the past participle. Neither subjunctive fragment supplies enough context to choose a full English conditional, wish, or other sentence. A response such as “they would have bound” invents tense and clause information absent from the task.
% mdparagraph:P020:end

\end{minipage}\par\addvspace{4pt}

\reserve{7\baselineskip}
% mdparagraph:P021:begin
\section*{P. Participles and passive clauses}
% mdparagraph:P021:end

% mdparagraph:P022:begin
Participles used adjectivally follow the assigned strong or weak agreement pattern. A verbal predicate permits an uninflected past participle or a strong-masculine-like ending; full attributive agreement is not compulsory there. Some spellings can occur in both constructions. [2]
% mdparagraph:P022:end

\par\begin{minipage}{\linewidth}\raggedright
% item: NEW-U07-P01
% key: KEY-NEW-U07-P01
% mdparagraph:P023:begin
\textbf{P01. Present participles.}
% mdparagraph:P023:end

% mdtable:T09:begin
\par\begingroup
\fontsize{10.5}{13.2}\selectfont
\setlength{\LTleft}{0pt}\setlength{\LTright}{\fill}
\noindent
\begin{tabular}{@{}>{\raggedright\arraybackslash}p{29.74mm}>{\raggedright\arraybackslash}p{30.49mm}>{\raggedright\arraybackslash}p{58.15mm}@{}}
\toprule
\textbf{Infinitive} & \textbf{Retained stem} & \textbf{Uninflected present participle} \\
\midrule
bindan & bind- & bindende \\
nęrian & nęri- & nęriende \\
lufian & lufi- & lufiende \\
cweðan & cweð- & cweðende \\
\bottomrule
\end{tabular}
\par\endgroup

% mdtable:T09:end
% mdparagraph:P024:begin
\textbf{Target:} add \textbf{-ende} to the present stem. In \textbf{nęrian} and \textbf{lufian}, remove \textbf{-an}, retaining the \textbf{i}: these are not forms made by removing all of \textbf{-ian}. Nor should an indicative third-singular stem such as \textbf{cwið-} replace \textbf{cweð-}.
% mdparagraph:P024:end

\end{minipage}\par\addvspace{4pt}

\par\begin{minipage}{\linewidth}\raggedright
% item: NEW-U07-P02
% key: KEY-NEW-U07-P02
% mdparagraph:P025:begin
\textbf{P02. Past participles.}
% mdparagraph:P025:end

% mdtable:T10:begin
\par\begingroup
\fontsize{10.5}{13.2}\selectfont
\setlength{\LTleft}{0pt}\setlength{\LTright}{\fill}
\noindent
\begin{tabular}{@{}>{\raggedright\arraybackslash}p{26.19mm}>{\raggedright\arraybackslash}p{42.69mm}>{\raggedright\arraybackslash}p{49.5mm}@{}}
\toprule
\textbf{Infinitive} & \textbf{Selected past participle} & \textbf{Formation} \\
\midrule
bindan & bunden & strong stem bund- with -en \\
beran & boren & strong stem bor- with -en \\
hīeran & hīered & weak dental formation \\
gehālgian & gehālgod & weak-II -od; retain lexical ge- \\
\bottomrule
\end{tabular}
\par\endgroup

% mdtable:T10:end
% mdparagraph:P026:begin
\textbf{Target:} reproduce the selected participles, including their vowels and prefixes. \textbf{Bindan} and \textbf{beran} are strong: adding a weak dental ending to \textbf{bind-} or \textbf{ber-} does not produce their past participles. \textbf{Hīered} retains the vowel before \textbf{d} in this selected citation form; its preterite \textbf{hīerde} is not interchangeable with it. An attested prefixed variant may be valid language, but the requested forms here are \textbf{bunden, boren, hīered, gehālgod}.
% mdparagraph:P026:end

\end{minipage}\par\addvspace{4pt}

\par\begin{minipage}{\linewidth}\raggedright
% item: NEW-U07-P03
% key: KEY-NEW-U07-P03
% mdparagraph:P027:begin
\textbf{P03. Agreement of bindende.}
% mdparagraph:P027:end

% mdtable:T11:begin
\par\begingroup
\fontsize{10.5}{13.2}\selectfont
\setlength{\LTleft}{0pt}\setlength{\LTright}{\fill}
\noindent
\begin{tabular}{@{}>{\raggedright\arraybackslash}p{43.2mm}>{\raggedright\arraybackslash}p{26.98mm}>{\raggedright\arraybackslash}p{48.19mm}@{}}
\toprule
\textbf{Part} & \textbf{Answer} & \textbf{Analysis} \\
\midrule
(a) a binding man, nominative & bindende mann & strong masculine nominative singular \\
(b) the binding man, nominative & se bindenda mann & weak masculine nominative singular \\
(c) dative feminine singular & bindendre & strong \\
(d) genitive plural, definite & bindendra & weak \\
\bottomrule
\end{tabular}
\par\endgroup

% mdtable:T11:end
% mdparagraph:P028:begin
\textbf{Target:} the determiner and the required case control the choice. There is no indefinite article to translate “a” in (a). The definite \textbf{se} in (b) selects weak \textbf{bindenda}, not strong \textbf{bindende}. Parts (c) and (d) ask only for the participle; adding an unfamiliar noun is unnecessary. \textbf{Bindendra} also has a strong genitive-plural use, but the stated definite phrase selects its weak use here.
% mdparagraph:P028:end

\end{minipage}\par\addvspace{4pt}

\par\begin{minipage}{\linewidth}\raggedright
% item: NEW-U07-P04
% key: KEY-NEW-U07-P04
% mdparagraph:P029:begin
\textbf{P04. Agreement of bunden.}
% mdparagraph:P029:end

% mdtable:T12:begin
\par\begingroup
\fontsize{10.5}{13.2}\selectfont
\setlength{\LTleft}{0pt}\setlength{\LTright}{\fill}
\noindent
\begin{tabular}{@{}>{\raggedright\arraybackslash}p{62.59mm}>{\raggedright\arraybackslash}p{32.68mm}>{\raggedright\arraybackslash}p{23.1mm}@{}}
\toprule
\textbf{Part} & \textbf{Participial answer} & \textbf{Analysis} \\
\midrule
(a) a bound animal, accusative neuter singular & bunden & strong \\
(b) the bound animal, accusative neuter singular & bundene & weak \\
(c) of bound birds, genitive plural & bundenra & strong \\
(d) to the bound birds, dative plural & bundenum & weak \\
\bottomrule
\end{tabular}
\par\endgroup

% mdtable:T12:end
% mdparagraph:P030:begin
\textbf{Target:} adjectival agreement. The fuller model preserves \textbf{bunden-}; the accepted shorter alternatives are \textbf{bundne} in (b) and \textbf{bundnum} in (d). Examples of the first two phrases are \textbf{bunden dēor} and \textbf{ðæt bundene dēor}. In (c), “of” requests a genitive, not a newly introduced prepositional phrase. An uninflected verbal-predicate rule cannot replace the requested attributive forms.
% mdparagraph:P030:end

\end{minipage}\par\addvspace{4pt}

\par\begin{minipage}{\linewidth}\raggedright
% item: NEW-U07-P05
% key: KEY-NEW-U07-P05
% mdparagraph:P031:begin
\textbf{P05. Hē bint ðæt dēor.}
% mdparagraph:P031:end

% mdparagraph:P032:begin
\textbf{Model:} \textbf{Ðæt dēor is bunden.} — “The animal is bound.”
% mdparagraph:P032:end

% mdparagraph:P033:begin
\textbf{Target:} present passive with \textbf{wesan}. In the active, \textbf{hē} is the subject and \textbf{ðæt dēor} the accusative patient. In the passive, \textbf{ðæt dēor} is the nominative subject and \textbf{is} agrees with its singular number. The neuter nominative and accusative happen to look alike; their syntactic functions have changed. The agent is omitted as requested. A form of \textbf{weorðan} could make a passive, but would miss this task’s chosen auxiliary. [3]
% mdparagraph:P033:end

\end{minipage}\par\addvspace{4pt}

\par\begin{minipage}{\linewidth}\raggedright
% item: NEW-U07-P06
% key: KEY-NEW-U07-P06
% mdparagraph:P034:begin
\textbf{P06. Hīe bundon ðā dēor.}
% mdparagraph:P034:end

% mdparagraph:P035:begin
\textbf{Model:} \textbf{Ðā dēor wurdon bunden.} — “The animals were bound,” describing the binding event.
% mdparagraph:P035:end

% mdparagraph:P036:begin
\textbf{Target:} preterite passive with \textbf{weorðan}. \textbf{Wurdon} is third plural preterite indicative, agreeing with the new nominative subject \textbf{ðā dēor}. The active subject \textbf{hīe} is omitted. \textbf{Bundene} is also defensible as the predicate participle. \textbf{Wǣron bunden} can be good Old English, but does not satisfy the express request for \textbf{weorðan}. Do not replace plural \textbf{wurdon} with singular \textbf{wearð} merely because \textbf{dēor} has a zero plural ending.
% mdparagraph:P036:end

\end{minipage}\par\addvspace{4pt}

\par\begin{minipage}{\linewidth}\raggedright
% item: NEW-U07-P07
% key: KEY-NEW-U07-P07
% mdparagraph:P037:begin
\textbf{P07. Ðā dēor sindon bunden /\allowbreak{} Ðā dēor sindon bundene.}
% mdparagraph:P037:end

% mdparagraph:P038:begin
Both can mean \textbf{“The animals are bound.”} \textbf{Sindon} is the finite plural auxiliary. \textbf{Bunden} is uninflected in the first verbal predicate; \textbf{bundene} has the permitted strong-masculine-like plural ending in the second. Neither requires the neuter subject to become masculine.
% mdparagraph:P038:end

% mdparagraph:P039:begin
\textbf{Target:} distinguish the construction from attributive agreement. The selected strong attributive neuter nominative/\allowbreak{}accusative plural is \textbf{bundenu}, as in \textbf{bundenu dēor}, “bound animals”; the adjective model also accepts \textbf{bunden} and \textbf{bundene} in those cells. Thus neither spelling is categorically forbidden attributively. The difference is that an attributive participle occupies an agreeing adjective cell, whereas a verbal predicate permits the predicate patterns just described. A spelling alone cannot prove which construction is present. [2]
% mdparagraph:P039:end

\end{minipage}\par\addvspace{4pt}

\par\begin{minipage}{\linewidth}\raggedright
% item: NEW-U07-P08
% key: KEY-NEW-U07-P08
% mdparagraph:P040:begin
\textbf{P08. Bēo geblētsod /\allowbreak{} Bēon ðā fuglas gemanigfielde.}
% mdparagraph:P040:end

% mdtable:T13:begin
\par\begingroup
\fontsize{10.5}{13.2}\selectfont
\setlength{\LTleft}{0pt}\setlength{\LTright}{\fill}
\noindent
\begin{tabular}{@{}>{\raggedright\arraybackslash}p{23.53mm}>{\raggedright\arraybackslash}p{73.7mm}>{\raggedright\arraybackslash}p{21.15mm}@{}}
\toprule
\textbf{Utterance} & \textbf{Analysis} & \textbf{Translation} \\
\midrule
Bēo geblētsod. & bēo: second singular imperative in the stated command; geblētsod: past participle in its verbal predicate & Be blessed. \\
Bēon ðā fuglas gemanigfielde. & bēon: third plural present subjunctive; ðā fuglas: nominative subject; gemanigfielde: past participle with plural predicate ending & Let the birds be multiplied. \\
\bottomrule
\end{tabular}
\par\endgroup

% mdtable:T13:end
% mdparagraph:P041:begin
\textbf{Target:} a command to one male addressee versus a wish about plural birds. The first \textbf{bēo} is also formally compatible with a singular present subjunctive, so the expressly stated command setting resolves its use. \textbf{Geblētsod} is compatible with the male singular addressee; its spelling also serves as the uninflected predicate base. In the second, \textbf{bēon} is finite, not an infinitive. The auxiliary supplies mood, person, and number; the participle does not independently express those finite categories.
% mdparagraph:P041:end

\end{minipage}\par\addvspace{4pt}

\reserve{7\baselineskip}
% mdparagraph:P042:begin
\section*{M. Mood in clauses}
% mdparagraph:P042:end

% mdparagraph:P043:begin
The translations in this section distinguish an assertion, a command, a wish, and a conditional setting. Neither English “may/\allowbreak{}would” nor the Old English conjunction \textbf{gif} is a mechanical instruction to use a particular mood. [4]
% mdparagraph:P043:end

\par\begin{minipage}{\linewidth}\raggedright
% item: NEW-U07-M01
% key: KEY-NEW-U07-M01
% mdparagraph:P044:begin
\textbf{M01. Statement and wish.}
% mdparagraph:P044:end

% mdparagraph:P045:begin
\textbf{Hīe swimmað on ðǣm sǣ.} — “They swim in the sea.” \textbf{Swimmað} is present indicative, third plural: the speaker asserts the activity.
% mdparagraph:P045:end

% mdparagraph:P046:begin
\textbf{Swimmen hīe on ðǣm sǣ.} — “Let them swim in the sea” or “May they swim in the sea.” \textbf{Swimmen} is present subjunctive, third plural: the speaker expresses a wish or exhortation about them.
% mdparagraph:P046:end

% mdparagraph:P047:begin
\textbf{Target:} distinguish the two clause forces. Both retain \textbf{on} with dative \textbf{ðǣm sǣ} for the location. The selected productive \textbf{sǣ} is masculine. The first sentence does not itself say they are able or permitted to swim; the second does not assert that they actually do so.
% mdparagraph:P047:end

\end{minipage}\par\addvspace{4pt}

\par\begin{minipage}{\linewidth}\raggedright
% item: NEW-U07-M02
% key: KEY-NEW-U07-M02
% mdparagraph:P048:begin
\textbf{M02. Gāð /\allowbreak{} Gān wē.}
% mdparagraph:P048:end

% mdparagraph:P049:begin
\textbf{Gāð.} — “Go,” addressed to more than one person: second plural imperative.
% mdparagraph:P049:end

% mdparagraph:P050:begin
\textbf{Gān wē.} — “Let us go”: first plural present subjunctive used as an exhortation, including the speaker in \textbf{wē}.
% mdparagraph:P050:end

% mdparagraph:P051:begin
\textbf{Target:} person and inclusion, as well as ending. Both urge an action, but the second is not a first-person “imperative form” in the selected morphological system. \textbf{Gān} can also be an infinitive elsewhere; here its subject and use make it finite.
% mdparagraph:P051:end

\end{minipage}\par\addvspace{4pt}

\par\begin{minipage}{\linewidth}\raggedright
% item: NEW-U07-M03
% key: KEY-NEW-U07-M03
% mdparagraph:P052:begin
\textbf{M03. Capability, permission, wish.}
% mdparagraph:P052:end

% mdtable:T14:begin
\par\begingroup
\fontsize{10.5}{13.2}\selectfont
\setlength{\LTleft}{0pt}\setlength{\LTright}{\fill}
\noindent
\begin{tabular}{@{}>{\raggedright\arraybackslash}p{30.2mm}>{\raggedright\arraybackslash}p{90.99mm}@{}}
\toprule
\textbf{Old English} & \textbf{Translation and force} \\
\midrule
Hē mæg swimman. & He can swim: capability. \\
Hē mōt swimman. & He may/\allowbreak{}is permitted to swim: permission. \\
Swimme hē. & Let him swim/\allowbreak{}May he swim: the speaker’s wish or exhortation. \\
\bottomrule
\end{tabular}
\par\endgroup

% mdtable:T14:end
% mdparagraph:P053:begin
\textbf{Target:} distinguish lexical modal meaning from inflectional mood. \textbf{Mæg} and \textbf{mōt} are third singular present indicative, each with the infinitive \textbf{swimman}. \textbf{Swimme} is third singular present subjunctive. English “may” can translate either permission or a wish; explain the intended sense rather than using that word as a mood test.
% mdparagraph:P053:end

\end{minipage}\par\addvspace{4pt}

\par\begin{minipage}{\linewidth}\raggedright
% item: NEW-U07-M04
% key: KEY-NEW-U07-M04
% mdparagraph:P054:begin
\textbf{M04. An open condition.}
% mdparagraph:P054:end

% mdparagraph:P055:begin
\textbf{Gif hē hæfð męte, hē sęlð ðǣm fuglum męte.} — “If he has food, he gives the birds food.” “If he has food, he will give the birds food” can also convey the open prospective setting; it does not turn the Old English verb into a separately inflected future.
% mdparagraph:P055:end

% mdparagraph:P056:begin
\textbf{Hæfð} and \textbf{sęlð} are both present indicative, third singular. \textbf{Gif} introduces the condition; the second clause states what follows in that circumstance. \textbf{Męte} is accusative in each clause; \textbf{ðǣm fuglum} is the dative gift-recipient in the second.
% mdparagraph:P056:end

% mdparagraph:P057:begin
\textbf{Target:} \textbf{gif} does not itself require subjunctive morphology. The task presents a real possibility without the specially tentative setting of M05. Calling \textbf{hæfð} subjunctive because it follows “if” misses the visible indicative form.
% mdparagraph:P057:end

\end{minipage}\par\addvspace{4pt}

\par\begin{minipage}{\linewidth}\raggedright
% item: NEW-U07-M05
% key: KEY-NEW-U07-M05
% mdparagraph:P058:begin
\textbf{M05. A tentative request.}
% mdparagraph:P058:end

% mdparagraph:P059:begin
\textbf{Gif ðū mē helpe, wē magon ðā dēor bindan.} — “If you help me, we can bind the animals.” “If you should help me” can make the tentativeness more explicit in English.
% mdparagraph:P059:end

% mdparagraph:P060:begin
\textbf{Helpe} is present subjunctive, second singular, agreeing with \textbf{ðū}; \textbf{mē} is the dative governed by \textbf{helpan}. \textbf{Magon} is present indicative, first plural, agreeing with \textbf{wē}, and \textbf{bindan} is its infinitive; \textbf{ðā dēor} is accusative plural.
% mdparagraph:P060:end

% mdparagraph:P061:begin
\textbf{Target:} the condition presents the desired help tentatively; the main clause claims ability conditional on that help. A subjunctive conditional clause does not automatically make its main clause subjunctive. Nor is \textbf{helpe} first singular indicative here merely because that form has the same spelling: the subject is \textbf{ðū}.
% mdparagraph:P061:end

\end{minipage}\par\addvspace{4pt}

\par\begin{minipage}{\linewidth}\raggedright
% item: NEW-U07-M06
% key: KEY-NEW-U07-M06
% mdparagraph:P062:begin
\textbf{M06. Supplied unreal, nonpast setting.}
% mdparagraph:P062:end

% mdparagraph:P063:begin
\textbf{Gif hē cōme, ic him hulpe.} — “If he came, I would help him.” The task explicitly supplies the setting that he is not coming.
% mdparagraph:P063:end

% mdparagraph:P064:begin
\textbf{Cōme} is preterite subjunctive, third singular, of \textbf{cuman}. \textbf{Hulpe} is preterite subjunctive, first singular, of \textbf{helpan}; \textbf{him} is its dative object. The morphological label “preterite” does not force past time: in this supplied unreal setting these forms refer to a nonpast possibility contrary to the setting.
% mdparagraph:P064:end

% mdparagraph:P065:begin
\textbf{Target:} recognise and explain the provided example. The answer need not create another unreal conditional or infer this reading without the supplied context. “He came and I helped him” would incorrectly assert the events as completed facts.
% mdparagraph:P065:end

\end{minipage}\par\addvspace{4pt}

\par\begin{minipage}{\linewidth}\raggedright
% item: NEW-U07-M07
% key: KEY-NEW-U07-M07
% mdparagraph:P066:begin
\textbf{M07. Hē wolde faran.}
% mdparagraph:P066:end

% mdparagraph:P067:begin
\textbf{Model translation:} “He wanted to go” or “He wished to go.” The narrator explicitly asserts a past wish, so \textbf{wolde} is third singular preterite indicative of \textbf{willan}; \textbf{faran} is its infinitive.
% mdparagraph:P067:end

% mdparagraph:P068:begin
\textbf{Target:} recognise syncretism without discarding the stated meaning. \textbf{Wolde} can also be a singular preterite subjunctive. The spelling alone cannot decide between them, and neither spelling nor mood automatically requires the English auxiliary “would.” “He would go” is acceptable only if its context makes the intended past volition clear; it is not a uniquely required translation.
% mdparagraph:P068:end

\end{minipage}\par\addvspace{4pt}

\par\begin{minipage}{\linewidth}\raggedright
% item: NEW-U07-M08
% key: KEY-NEW-U07-M08
% mdparagraph:P069:begin
\textbf{M08. Make the wishes.}
% mdparagraph:P069:end

% mdparagraph:P070:begin
\textbf{Hē leofað → Libbe hē.} — “May he live.” \textbf{Libbe} is third singular present subjunctive, with the selected \textbf{libb-} stem.
% mdparagraph:P070:end

% mdparagraph:P071:begin
\textbf{Wē helpað ðǣm męnn → Helpen wē ðǣm męnn.} — “Let us help the man.” \textbf{Helpen} is first plural present subjunctive. \textbf{Ðǣm męnn} remains dative singular, as required by \textbf{helpan}.
% mdparagraph:P071:end

% mdparagraph:P072:begin
\textbf{Target:} change the clause force and verb form while preserving reference and government. \textbf{Hē libbe} and \textbf{Wē helpen ðǣm męnn} are defensible orders in an appropriate wish context. \textbf{Help} or \textbf{helpað} would instead address a second person; changing the object to accusative \textbf{ðone mann} would lose the taught frame.
% mdparagraph:P072:end

\end{minipage}\par\addvspace{4pt}

\reserve{7\baselineskip}
% mdparagraph:P073:begin
\section*{C. Short composition}
% mdparagraph:P073:end

\par\begin{minipage}{\linewidth}\raggedright
% item: NEW-U07-C01
% key: KEY-NEW-U07-C01
% mdparagraph:P074:begin
\textbf{C01. “May the birds live.”}
% mdparagraph:P074:end

% mdparagraph:P075:begin
\textbf{Model:} \textbf{Libben ðā fuglas.}
% mdparagraph:P075:end

% mdparagraph:P076:begin
\textbf{Target:} third plural present subjunctive of \textbf{libban}, with nominative plural \textbf{ðā fuglas}. The essential meaning is a wish for the birds’ life, not an assertion that they live. \textbf{Ðā fuglas libben} is possible with the same intended force. \textbf{Leofað} is the selected third-singular indicative; \textbf{libbað} is the plural indicative or imperative. Neither supplies the requested plural wish.
% mdparagraph:P076:end

\end{minipage}\par\addvspace{4pt}

\par\begin{minipage}{\linewidth}\raggedright
% item: NEW-U07-C02
% key: KEY-NEW-U07-C02
% mdparagraph:P077:begin
\textbf{C02. “Let us help the companions.”}
% mdparagraph:P077:end

% mdparagraph:P078:begin
\textbf{Model:} \textbf{Helpen wē ðǣm gefērum.}
% mdparagraph:P078:end

% mdparagraph:P079:begin
\textbf{Target:} first plural exhortation and the dative object of \textbf{helpan}. The speaker belongs to \textbf{wē}, and \textbf{ðǣm gefērum} identifies the companions receiving help. Reordering the phrase can be valid. \textbf{Helpað} addresses others; accusative \textbf{ðā gefēran} would not satisfy this verb’s government.
% mdparagraph:P079:end

\end{minipage}\par\addvspace{4pt}

\par\begin{minipage}{\linewidth}\raggedright
% item: NEW-U07-C03
% key: KEY-NEW-U07-C03
% mdparagraph:P080:begin
\textbf{C03. Tentative: “If you (sing.) have food, give it to the birds.”}
% mdparagraph:P080:end

% mdparagraph:P081:begin
\textbf{Model:} \textbf{Gif ðū hæbbe męte, sęle hine ðǣm fuglum.}
% mdparagraph:P081:end

% mdparagraph:P082:begin
\textbf{Target:} \textbf{hæbbe}, present subjunctive second singular, in the tentative condition; \textbf{sęle}, second singular imperative, in the command. \textbf{Hine} is accusative masculine singular because it refers to masculine \textbf{męte}, even though English calls food “it.” \textbf{Ðǣm fuglum} is the dative recipient.
% mdparagraph:P082:end

% mdparagraph:P083:begin
\textbf{Gief hine ðǣm fuglum} is a defensible alternative using \textbf{giefan}. Indicative \textbf{hæfst} can form a valid open condition, but does not meet the expressly tentative target. \textbf{Hit} would miss the selected antecedent’s grammatical gender; indicative \textbf{sęlð} would miss the singular command.
% mdparagraph:P083:end

\end{minipage}\par\addvspace{4pt}

\par\begin{minipage}{\linewidth}\raggedright
% item: NEW-U07-C04
% key: KEY-NEW-U07-C04
% mdparagraph:P084:begin
\textbf{C04. Command to one male addressee: “Be blessed.”}
% mdparagraph:P084:end

% mdparagraph:P085:begin
\textbf{Model:} \textbf{Bēo geblētsod.}
% mdparagraph:P085:end

% mdparagraph:P086:begin
\textbf{Target:} singular passive command, using the selected participle of \textbf{geblētsian}. The stated setting selects imperative \textbf{bēo}, although it is formally also a singular subjunctive. \textbf{Bēoð} would address more than one person. The supplied unprefixed \textbf{blētsode} in the source reading does not establish an independently learned unprefixed production paradigm.
% mdparagraph:P086:end

\end{minipage}\par\addvspace{4pt}

\par\begin{minipage}{\linewidth}\raggedright
% item: NEW-U07-C05
% key: KEY-NEW-U07-C05
% mdparagraph:P087:begin
\textbf{C05. Wish: “Let the work be accomplished.”}
% mdparagraph:P087:end

% mdparagraph:P088:begin
\textbf{Model:} \textbf{Bēo ðæt weorc gefręmed.}
% mdparagraph:P088:end

% mdparagraph:P089:begin
\textbf{Target:} third singular present subjunctive passive. \textbf{Ðæt weorc} is nominative neuter singular; \textbf{gefręmed} is the selected past participle of \textbf{gefręmman}. \textbf{Sīe ðæt weorc gefręmed} is equally suitable. \textbf{Is} would assert the state; an active form of \textbf{gefręmman} would change the relation between the work and the accomplishing.
% mdparagraph:P089:end

\end{minipage}\par\addvspace{4pt}

\par\begin{minipage}{\linewidth}\raggedright
% item: NEW-U07-C06
% key: KEY-NEW-U07-C06
% mdparagraph:P090:begin
\textbf{C06. “Every bird is above the earth; all trees bear fruit.”}
% mdparagraph:P090:end

% mdparagraph:P091:begin
\textbf{Model:} \textbf{Ǣlc fugol is bufan ðǣre eorðan; eall trēowu berað wæstm.}
% mdparagraph:P091:end

% mdparagraph:P092:begin
\textbf{Target:} singular \textbf{ǣlc fugol}, plural \textbf{eall trēowu}, and dative after \textbf{bufan}. \textbf{Is} agrees with singular \textbf{fugol}; \textbf{berað} agrees with plural \textbf{trēowu} and illustrates the new produce-bearing sense of \textbf{beran}. \textbf{Wæstm} is accusative singular used for fruit collectively; \textbf{wæstmas} can express individual fruits. \textbf{Wyrceað wæstm}, “yield fruit,” is also a valid alternative. Do not make the subject plural merely because “every” ranges over many birds.
% mdparagraph:P092:end

\end{minipage}\par\addvspace{4pt}

\reserve{7\baselineskip}
% mdparagraph:P093:begin
\section*{Cook I.4.18 and I.6–7}
% mdparagraph:P093:end

% mdparagraph:P094:begin
The English prompts below preserve Cook’s wording and original identifiers. The Old English models are new teaching compositions, not transcriptions of an answer supplied by Cook. Word order can vary where the cases and clause force remain clear. The references to number in Cook’s parentheses are binding; an unqualified English “you” permits a suitable singular, plural, or explicitly dual referent. [5]
% mdparagraph:P094:end

\par\begin{minipage}{\linewidth}\raggedright
% item: EX-I-04-18
% key: KEY-EX-I-04-18
% mdparagraph:P095:begin
\textbf{I.4.18. Increase (sing.) and be multiplied.}
% mdparagraph:P095:end

% mdparagraph:P096:begin
\textbf{Model:} \textbf{Weax and bēo gemanigfield.}
% mdparagraph:P096:end

% mdparagraph:P097:begin
\textbf{Target:} two commands to one addressee, the second passive. \textbf{Weax} is the singular imperative of \textbf{weaxan}; \textbf{bēo} is singular imperative in this command context. \textbf{Gemanigfield} is the uninflected predicate participle and does not impose an unstated gender on the addressee. \textbf{Weaxað} and \textbf{bēoð} would wrongly address a plural group. Cook’s note to Exercise I.4 is “[NOTE.—For adjectives preceded by a demonstrative, see 55.]”; no separate note is attached to item 18.
% mdparagraph:P097:end

\end{minipage}\par\addvspace{4pt}

\reserve{7\baselineskip}
% mdparagraph:P098:begin
\subsection*{Cook I.6}
% mdparagraph:P098:end

\par\begin{minipage}{\linewidth}\raggedright
% item: EX-I-06-01
% key: KEY-EX-I-06-01
% mdparagraph:P099:begin
\textbf{I.6.1. The animal lives.}
% mdparagraph:P099:end

% mdparagraph:P100:begin
\textbf{Model:} \textbf{Ðæt dēor leofað.}
% mdparagraph:P100:end

% mdparagraph:P101:begin
\textbf{Target:} singular subject and weak-III present indicative. \textbf{Ðæt dēor} is nominative neuter singular; \textbf{leofað} is third singular. Although \textbf{dēor} also has a zero plural form, the prompt and singular demonstrative determine the number here. \textbf{Libbe} would express a different mood in a third-person clause.
% mdparagraph:P101:end

\end{minipage}\par\addvspace{4pt}

\par\begin{minipage}{\linewidth}\raggedright
% item: EX-I-06-02
% key: KEY-EX-I-06-02
% mdparagraph:P102:begin
\textbf{I.6.2. The bird flies.}
% mdparagraph:P102:end

% mdparagraph:P103:begin
\textbf{Model:} \textbf{Se fugol flīehð.}
% mdparagraph:P103:end

% mdparagraph:P104:begin
\textbf{Target:} nominative masculine singular and the selected third-singular present stem of \textbf{flēogan}. Retain the vowel and consonant alternation in \textbf{flīehð}. The infinitive \textbf{flēogan} cannot serve as the sole finite verb of this assertion.
% mdparagraph:P104:end

\end{minipage}\par\addvspace{4pt}

\par\begin{minipage}{\linewidth}\raggedright
% item: EX-I-06-03
% key: KEY-EX-I-06-03
% mdparagraph:P105:begin
\textbf{I.6.3. The tree grows.}
% mdparagraph:P105:end

% mdparagraph:P106:begin
\textbf{Model:} \textbf{Ðæt trēo wiext.}
% mdparagraph:P106:end

% mdparagraph:P107:begin
\textbf{Target:} nominative neuter singular with third singular indicative \textbf{wiext}. Present subjunctive \textbf{weaxe} would make a wish in a suitable setting; it would not reproduce this plain assertion. Preserve the selected stem alternation rather than adding an ending mechanically to \textbf{weax-}.
% mdparagraph:P107:end

\end{minipage}\par\addvspace{4pt}

\par\begin{minipage}{\linewidth}\raggedright
% item: EX-I-06-04
% key: KEY-EX-I-06-04
% mdparagraph:P108:begin
\textbf{I.6.4. The man rules over the cattle of the earth.}
% mdparagraph:P108:end

% mdparagraph:P109:begin
\textbf{Model:} \textbf{Se mann gīemð ðǣre eorðan nīetena.}
% mdparagraph:P109:end

% mdparagraph:P110:begin
\textbf{Target:} \textbf{gīeman} governs a genitive in this sense. \textbf{Nīetena} is genitive plural, the cattle ruled over; \textbf{ðǣre eorðan} is a dependent genitive singular, “of the earth,” identifying those cattle. The verb agrees with singular \textbf{se mann}. An order such as \textbf{nīetena ðǣre eorðan} can preserve the relationships. An accusative direct object after \textbf{gīemð} would miss its taught government. English “over” does not require an Old English preposition when the chosen verb already governs the relation.
% mdparagraph:P110:end

\end{minipage}\par\addvspace{4pt}

\par\begin{minipage}{\linewidth}\raggedright
% item: EX-I-06-05
% key: KEY-EX-I-06-05
% mdparagraph:P111:begin
\textbf{I.6.5. Let him be over the birds and (over) the beasts.}
% mdparagraph:P111:end

% mdparagraph:P112:begin
\textbf{Model:} \textbf{Bēo hē ofer ðā fuglas and ofer ðā dēor.}
% mdparagraph:P112:end

% mdparagraph:P113:begin
\textbf{Target:} third singular present subjunctive and \textbf{ofer} with accusative in the authority relation. \textbf{Sīe hē} is equally suitable. The second \textbf{ofer} may be omitted, as Cook’s parentheses indicate: \textbf{ofer ðā fuglas and ðā dēor}. Here “over” concerns dominion; it is not a request to apply a mechanical physical-location rule. A second-person imperative would change the person intended. [6]
% mdparagraph:P113:end

\end{minipage}\par\addvspace{4pt}

\par\begin{minipage}{\linewidth}\raggedright
% item: EX-I-06-06
% key: KEY-EX-I-06-06
% mdparagraph:P114:begin
\textbf{I.6.6. The stars are above the sea.}
% mdparagraph:P114:end

% mdparagraph:P115:begin
\textbf{Model:} \textbf{Ðā steorran sindon bufan ðǣm sǣ.}
% mdparagraph:P115:end

% mdparagraph:P116:begin
\textbf{Target:} plural subject, plural indicative copula, and dative after \textbf{bufan}. The course’s productive \textbf{sǣ} is masculine, hence \textbf{ðǣm}. The feminine genitive in the source reading is locally supplied and is not the productive paradigm requested here. Attested feminine uses of \textbf{sǣ} are not errors in Old English, but require their own appropriate agreement and evidence.
% mdparagraph:P116:end

\end{minipage}\par\addvspace{4pt}

\par\begin{minipage}{\linewidth}\raggedright
% item: EX-I-06-07
% key: KEY-EX-I-06-07
% mdparagraph:P117:begin
\textbf{I.6.7. It is for your food [to you for food].}
% mdparagraph:P117:end

% mdparagraph:P118:begin
\textbf{Model:} \textbf{Hit is ēow tō męte.}
% mdparagraph:P118:end

% mdparagraph:P119:begin
\textbf{Target:} dative recipient \textbf{ēow}, with \textbf{tō męte} expressing the use “for food.” The square-bracketed wording is Cook’s own help: the model literally says “It is to you for food.” \textbf{Ðē} is equally valid for one addressee; \textbf{inc} is possible for two. Cook does not specify number, so \textbf{ēow} is not a compulsory hidden plural target. \textbf{Męte} is dative after \textbf{tō} here, though other cases share its spelling.
% mdparagraph:P119:end

\end{minipage}\par\addvspace{4pt}

\par\begin{minipage}{\linewidth}\raggedright
% item: EX-I-06-08
% key: KEY-EX-I-06-08
% mdparagraph:P120:begin
\textbf{I.6.8. The herb bears seed.}
% mdparagraph:P120:end

% mdparagraph:P121:begin
\textbf{Model:} \textbf{Sēo wyrt birð sǣd.}
% mdparagraph:P121:end

% mdparagraph:P122:begin
\textbf{Target:} the produce-bearing sense of \textbf{beran}, with a singular subject and accusative object. \textbf{Sēo wyrt} is feminine nominative singular; \textbf{birð} is third singular present indicative. \textbf{Sǣd} is accusative, here collective seed. Translating the verb merely as physical carrying would miss the sense required by a growing herb. A suitable count-plural reading of seed can be defensible where the form and context support it.
% mdparagraph:P122:end

\end{minipage}\par\addvspace{4pt}

\par\begin{minipage}{\linewidth}\raggedright
% item: EX-I-06-09
% key: KEY-EX-I-06-09
% mdparagraph:P123:begin
\textbf{I.6.9. The tree yields fruit.}
% mdparagraph:P123:end

% mdparagraph:P124:begin
\textbf{Model:} \textbf{Ðæt trēo wyrcð wæstm.}
% mdparagraph:P124:end

% mdparagraph:P125:begin
\textbf{Target:} \textbf{wyrcean} in the taught yield/\allowbreak{}produce sense; \textbf{wyrcð} is third singular, and \textbf{wæstm} the accusative yield. \textbf{Birð wæstm} is a defensible lexical alternative expressing fruit-bearing. \textbf{Wæstmas} can refer to individual fruits; collective singular \textbf{wæstm} closely matches the model. Do not replace the finite verb with its infinitive or confuse singular \textbf{wyrcð} with plural \textbf{wyrceað}.
% mdparagraph:P125:end

\end{minipage}\par\addvspace{4pt}

\par\begin{minipage}{\linewidth}\raggedright
% item: EX-I-06-10
% key: KEY-EX-I-06-10
% mdparagraph:P126:begin
\textbf{I.6.10. The air is very good.}
% mdparagraph:P126:end

% mdparagraph:P127:begin
\textbf{Model:} \textbf{Ðæt lyft is swīðe gōd.}
% mdparagraph:P127:end

% mdparagraph:P128:begin
\textbf{Target:} the course’s selected neuter \textbf{lyft}, singular copula, and strong predicate adjective. \textbf{Gōd} is nominative neuter singular; \textbf{swīðe} is the invariant adverb “very.” The noun’s gender must come from the selected lexical entry, not an assumed modern-language gender or a different attested paradigm.
% mdparagraph:P128:end

\end{minipage}\par\addvspace{4pt}

\par\begin{minipage}{\linewidth}\raggedright
% item: EX-I-06-11
% key: KEY-EX-I-06-11
% mdparagraph:P129:begin
\textbf{I.6.11. The things are very good.}
% mdparagraph:P129:end

% mdparagraph:P130:begin
\textbf{Model:} \textbf{Ðā ðing sindon swīðe gōd.}
% mdparagraph:P130:end

% mdparagraph:P131:begin
\textbf{Target:} nominative neuter plural \textbf{ðing}, plural \textbf{sindon}, and the selected strong neuter plural \textbf{gōd}. The adjective’s zero ending does not make it singular. The demonstrative belongs to the subject noun phrase; it does not automatically make the separate predicate adjective weak.
% mdparagraph:P131:end

\end{minipage}\par\addvspace{4pt}

\par\begin{minipage}{\linewidth}\raggedright
% item: EX-I-06-12
% key: KEY-EX-I-06-12
% mdparagraph:P132:begin
\textbf{I.6.12. The birds were very good.}
% mdparagraph:P132:end

% mdparagraph:P133:begin
\textbf{Model:} \textbf{Ðā fuglas wǣron swīðe gōde.}
% mdparagraph:P133:end

% mdparagraph:P134:begin
\textbf{Target:} past plural \textbf{wǣron}, with strong nominative masculine plural \textbf{gōde}. Contrast neuter plural \textbf{gōd} in I.6.11. \textbf{Sindon} would lose the past time; \textbf{gōdan} would substitute a weak adjective form for the selected strong predicate agreement.
% mdparagraph:P134:end

\end{minipage}\par\addvspace{4pt}

\par\begin{minipage}{\linewidth}\raggedright
% item: EX-I-06-13
% key: KEY-EX-I-06-13
% mdparagraph:P135:begin
\textbf{I.6.13. The light of the stars was very good.}
% mdparagraph:P135:end

% mdparagraph:P136:begin
\textbf{Model:} \textbf{Ðāra steorrena lēoht wæs swīðe gōd.}
% mdparagraph:P136:end

% mdparagraph:P137:begin
\textbf{Target:} a dependent genitive plural within a singular subject phrase. \textbf{Ðāra steorrena} means “of the stars”; \textbf{lēoht} is the nominative neuter singular head, controlling \textbf{wæs} and \textbf{gōd}. An order with the genitive after \textbf{lēoht} can preserve the meaning. Making the verb plural because the stars are plural would attach agreement to the wrong noun.
% mdparagraph:P137:end

\end{minipage}\par\addvspace{4pt}

\par\begin{minipage}{\linewidth}\raggedright
% item: EX-I-06-14
% key: KEY-EX-I-06-14
% mdparagraph:P138:begin
\textbf{I.6.14. He gave them all trees.}
% mdparagraph:P138:end

% mdparagraph:P139:begin
\textbf{Model:} \textbf{Hē geaf him eall trēowu.}
% mdparagraph:P139:end

% mdparagraph:P140:begin
\textbf{Target:} past giving with dative recipient \textbf{him}, here plural “them,” and accusative plural gift \textbf{eall trēowu}. \textbf{Eall} agrees with neuter plural \textbf{trēowu}. \textbf{Hē sealde him eall trēowu} is a defensible alternative with \textbf{sęllan}. The spelling \textbf{him} also serves singular datives, but the English prompt fixes plural reference in this answer. Do not turn the trees into the recipients.
% mdparagraph:P140:end

\end{minipage}\par\addvspace{4pt}

\par\begin{minipage}{\linewidth}\raggedright
% item: EX-I-06-15
% key: KEY-EX-I-06-15
% mdparagraph:P141:begin
\textbf{I.6.15. I give you all things upon earth.}
% mdparagraph:P141:end

% mdparagraph:P142:begin
\textbf{Model:} \textbf{Ic giefe ēow eall ðing on eorðan.}
% mdparagraph:P142:end

% mdparagraph:P143:begin
\textbf{Target:} first singular present giving, dative recipient, accusative neuter plural gift, and the location “on earth.” \textbf{Ðē} may replace \textbf{ēow} for one addressee; \textbf{inc} is appropriate if two are intended. \textbf{Ic sęlle} is a defensible lexical alternative. Preserve “all things,” not merely one thing, despite the zero plural forms \textbf{eall ðing}. The original does not require plural “you.”
% mdparagraph:P143:end

\end{minipage}\par\addvspace{4pt}

\par\begin{minipage}{\linewidth}\raggedright
% item: EX-I-06-16
% key: KEY-EX-I-06-16
% mdparagraph:P144:begin
\textbf{I.6.16. I call the firmament Heaven.}
% mdparagraph:P144:end

% mdparagraph:P145:begin
\textbf{Model:} \textbf{Ic hāte ðā fæstnesse Heofon.}
% mdparagraph:P145:end

% mdparagraph:P146:begin
\textbf{Target:} naming with two accusatives: the entity named, \textbf{ðā fæstnesse}, and the name assigned, \textbf{Heofon}. \textbf{Hāte} is first singular present indicative. \textbf{Heofon} is capitalised as the assigned name. It is not a second nominative subject, nor is this sentence a passive “I am called.” [6]
% mdparagraph:P146:end

\end{minipage}\par\addvspace{4pt}

\reserve{7\baselineskip}
% mdparagraph:P147:begin
\subsection*{Cook I.7}
% mdparagraph:P147:end

\par\begin{minipage}{\linewidth}\raggedright
% item: EX-I-07-01
% key: KEY-EX-I-07-01
% mdparagraph:P148:begin
\textbf{I.7.1. He has dominion.}
% mdparagraph:P148:end

% mdparagraph:P149:begin
\textbf{Model:} \textbf{Hē rīcsað.}
% mdparagraph:P149:end

% mdparagraph:P150:begin
\textbf{Target:} \textbf{rīcsian} used absolutely, “exercise dominion, rule.” The English phrase need not be reproduced word for word with \textbf{habban}. A well-formed equivalent using “has dominion” can preserve the meaning, while the model practises this new verb. \textbf{Rīcsien} would be subjunctive plural and would not express the singular assertion.
% mdparagraph:P150:end

\end{minipage}\par\addvspace{4pt}

\par\begin{minipage}{\linewidth}\raggedright
% item: EX-I-07-02
% key: KEY-EX-I-07-02
% mdparagraph:P151:begin
\textbf{I.7.2. He blesses the earth.}
% mdparagraph:P151:end

% mdparagraph:P152:begin
\textbf{Model:} \textbf{Hē geblētsað ðā eorðan.}
% mdparagraph:P152:end

% mdparagraph:P153:begin
\textbf{Target:} third singular present indicative and an accusative object. The model uses previously taught \textbf{geblētsian}. \textbf{Ðā eorðan} is feminine accusative singular, not a plural phrase. The reading’s plain \textbf{blētsode} is separately supplied; its occurrence does not require removal of the prefix in independent composition.
% mdparagraph:P153:end

\end{minipage}\par\addvspace{4pt}

\par\begin{minipage}{\linewidth}\raggedright
% item: EX-I-07-03
% key: KEY-EX-I-07-03
% mdparagraph:P154:begin
\textbf{I.7.3. He says thus.}
% mdparagraph:P154:end

% mdparagraph:P155:begin
\textbf{Model:} \textbf{Hē sægð ðus.}
% mdparagraph:P155:end

% mdparagraph:P156:begin
\textbf{Target:} the selected third singular \textbf{sægð}, with adverb \textbf{ðus}, “thus.” Preserve the lexical present alternation of \textbf{sęcgan}. “Thus” is an adverb here; it does not introduce a relative clause or take an agreement ending.
% mdparagraph:P156:end

\end{minipage}\par\addvspace{4pt}

\par\begin{minipage}{\linewidth}\raggedright
% item: EX-I-07-04
% key: KEY-EX-I-07-04
% mdparagraph:P157:begin
\textbf{I.7.4. Now is the season.}
% mdparagraph:P157:end

% mdparagraph:P158:begin
\textbf{Model:} \textbf{Nū is sēo tīd.}
% mdparagraph:P158:end

% mdparagraph:P159:begin
\textbf{Target:} temporal \textbf{nū}, feminine nominative singular \textbf{sēo tīd}, and singular \textbf{is}. “Time” can convey the ordinary sense of \textbf{tīd}, but the given English “season” is also the required meaning here. \textbf{Sēo tīd is nū} can preserve the basic predication; no change to subjunctive mood is licensed merely by initial \textbf{nū}.
% mdparagraph:P159:end

\end{minipage}\par\addvspace{4pt}

\par\begin{minipage}{\linewidth}\raggedright
% item: EX-I-07-05
% key: KEY-EX-I-07-05
% mdparagraph:P160:begin
\textbf{I.7.5. It was so called.}
% mdparagraph:P160:end

% mdparagraph:P161:begin
\textbf{Model:} \textbf{Hit wæs swā hāten.}
% mdparagraph:P161:end

% mdparagraph:P162:begin
\textbf{Target:} past passive with \textbf{wesan}, participle \textbf{hāten}, and adverb \textbf{swā}, “so.” \textbf{Wæs} gives past singular predication; \textbf{hāten} does not itself carry finite tense or person. An attested \textbf{gehāten} variant can be defensible. A correct use of inherited passive \textbf{hātte} is not bad Old English, but that special paradigm is not this unit’s production target. Do not confuse active “it called” with “it was called.”
% mdparagraph:P162:end

\end{minipage}\par\addvspace{4pt}

\par\begin{minipage}{\linewidth}\raggedright
% item: EX-I-07-06
% key: KEY-EX-I-07-06
% mdparagraph:P163:begin
\textbf{I.7.6. Behold, God rules over the earth.}
% mdparagraph:P163:end

% mdparagraph:P164:begin
\textbf{Model:} \textbf{Efne, God gīemð ðǣre eorðan.}
% mdparagraph:P164:end

% mdparagraph:P165:begin
\textbf{Target:} attention-directing \textbf{efne}, explicit subject \textbf{God}, and genitive after \textbf{gīeman}. \textbf{Ðǣre eorðan} is genitive singular in this frame. English “behold” is rendered by the learned particle, without requiring a new visual-perception command. An accusative earth-phrase after \textbf{gīemð} would lose the selected government.
% mdparagraph:P165:end

\end{minipage}\par\addvspace{4pt}

\par\begin{minipage}{\linewidth}\raggedright
% item: EX-I-07-07
% key: KEY-EX-I-07-07
% mdparagraph:P166:begin
\textbf{I.7.7. The tree is in the midst of the grass.}
% mdparagraph:P166:end

% mdparagraph:P167:begin
\textbf{Model:} \textbf{Ðæt trēo is tōmiddes ðǣm gærse.}
% mdparagraph:P167:end

% mdparagraph:P168:begin
\textbf{Target:} \textbf{tōmiddes} with dative in the selected spatial frame. \textbf{Ðǣm gærse} is dative singular; \textbf{ðæt trēo} is the singular subject. The English word “of” is part of “in the midst of”; it does not require a genitive where the taught Old English preposition governs a dative.
% mdparagraph:P168:end

\end{minipage}\par\addvspace{4pt}

\par\begin{minipage}{\linewidth}\raggedright
% item: EX-I-07-08
% key: KEY-EX-I-07-08
% mdparagraph:P169:begin
\textbf{I.7.8. The earth is in the midst of the air.}
% mdparagraph:P169:end

% mdparagraph:P170:begin
\textbf{Model:} \textbf{Sēo eorðe is tōmiddes ðǣm lyfte.}
% mdparagraph:P170:end

% mdparagraph:P171:begin
\textbf{Target:} feminine nominative \textbf{sēo eorðe} as subject, and the selected neuter dative \textbf{ðǣm lyfte} after \textbf{tōmiddes}. The two nouns have different genders and different syntactic roles. Repeating nominative \textbf{eorðe} as an oblique form or selecting a feminine article for this course’s \textbf{lyft} would miss those distinctions.
% mdparagraph:P171:end

\end{minipage}\par\addvspace{4pt}

\par\begin{minipage}{\linewidth}\raggedright
% item: EX-I-07-09
% key: KEY-EX-I-07-09
% mdparagraph:P172:begin
\textbf{I.7.9. The waters are under the firmament.}
% mdparagraph:P172:end

% mdparagraph:P173:begin
\textbf{Model:} \textbf{Ðā wæteru sindon under ðǣre fæstnesse.}
% mdparagraph:P173:end

% mdparagraph:P174:begin
\textbf{Target:} plural subject and \textbf{under} with dative for this location. \textbf{Ðǣre fæstnesse} is feminine dative singular; \textbf{ðā wæteru} is neuter nominative plural. The location does not introduce a second subject, and the singular firmament does not control the verb’s number.
% mdparagraph:P174:end

\end{minipage}\par\addvspace{4pt}

\par\begin{minipage}{\linewidth}\raggedright
% item: EX-I-07-10
% key: KEY-EX-I-07-10
% mdparagraph:P175:begin
\textbf{I.7.10. Birds are the food of the animal.}
% mdparagraph:P175:end

% mdparagraph:P176:begin
\textbf{Model:} \textbf{Fuglas sindon ðæs dēores męte.}
% mdparagraph:P176:end

% mdparagraph:P177:begin
\textbf{Target:} plural \textbf{fuglas} as subject, plural \textbf{sindon}, and predicate noun \textbf{męte} with dependent genitive \textbf{ðæs dēores}. \textbf{Męte} remains singular: a predicate noun need not copy the subject’s number. The genitive phrase identifies whose food it is; it does not make the animal the subject. An article before \textbf{fuglas} would make the birds more specifically definite than Cook’s bare plural.
% mdparagraph:P177:end

\end{minipage}\par\addvspace{4pt}

\par\begin{minipage}{\linewidth}\raggedright
% item: EX-I-07-11
% key: KEY-EX-I-07-11
% mdparagraph:P178:begin
\textbf{I.7.11. Divide (plur.) the waters from the waters.}
% mdparagraph:P178:end

% mdparagraph:P179:begin
\textbf{Model:} \textbf{Tōtwǣmað ðā wæteru fram ðǣm wæterum.}
% mdparagraph:P179:end

% mdparagraph:P180:begin
\textbf{Target:} second plural imperative \textbf{tōtwǣmað}, accusative thing divided, and dative after \textbf{fram}. The first water-phrase is accusative plural; the second is dative plural. Although a plural indicative can share the ending \textbf{-að}, the command context determines its use here. \textbf{Tōtwǣm} would address only one person and miss Cook’s explicit number.
% mdparagraph:P180:end

\end{minipage}\par\addvspace{4pt}

\par\begin{minipage}{\linewidth}\raggedright
% item: EX-I-07-12
% key: KEY-EX-I-07-12
% mdparagraph:P181:begin
\textbf{I.7.12. God's dominion is over all creatures.}
% mdparagraph:P181:end

% mdparagraph:P182:begin
\textbf{Model:} \textbf{Godes geweald is ofer ealla gesceafta.}
% mdparagraph:P182:end

% mdparagraph:P183:begin
\textbf{Target:} singular \textbf{geweald}, dependent genitive \textbf{Godes}, and the authority use of \textbf{ofer} with accusative. \textbf{Ealla gesceafta} is feminine accusative plural. \textbf{Ealle gesceafta} is also defensible and is attested in Cook’s Creation text. Neither plural creatures nor the genitive possessor replaces \textbf{geweald} as the head controlling \textbf{is}. The form does not demand a plural “dominions” sense. [6]
% mdparagraph:P183:end

\end{minipage}\par\addvspace{4pt}

\par\begin{minipage}{\linewidth}\raggedright
% item: EX-I-07-13
% key: KEY-EX-I-07-13
% mdparagraph:P184:begin
\textbf{I.7.13. God hallows each day.}
% mdparagraph:P184:end

% mdparagraph:P185:begin
\textbf{Model:} \textbf{God gehālgað ǣlcne dæg.}
% mdparagraph:P185:end

% mdparagraph:P186:begin
\textbf{Target:} keep \textbf{God} as the explicit subject, third singular present \textbf{gehālgað}, and singular accusative \textbf{ǣlcne dæg}. The ending \textbf{-ne} agrees with masculine accusative singular \textbf{dæg}. “Each” distributes the action over days while retaining a singular noun phrase. Replacing \textbf{God} with an unexplained “he” or turning “each day” into plural “all days” loses part of the original prompt.
% mdparagraph:P186:end

\end{minipage}\par\addvspace{4pt}

\par\begin{minipage}{\linewidth}\raggedright
% item: EX-I-07-14
% key: KEY-EX-I-07-14
% mdparagraph:P187:begin
\textbf{I.7.14. Let the birds increase and be multiplied.}
% mdparagraph:P187:end

% mdparagraph:P188:begin
\textbf{Model:} \textbf{Weaxen ðā fuglas and bēon gemanigfielde.}
% mdparagraph:P188:end

% mdparagraph:P189:begin
\textbf{Target:} coordinated third plural wishes, the first active and the second passive. \textbf{Ðā fuglas} supplies the subject of both; \textbf{gemanigfielde} has the permitted plural predicate ending. \textbf{Sīen} can replace \textbf{bēon}, and an uninflected \textbf{gemanigfield} can also serve in the verbal predicate. \textbf{Weaxað … bēoð} would instead make commands addressed to a plural group; Cook’s wording requests a wish about the birds.
% mdparagraph:P189:end

\end{minipage}\par\addvspace{4pt}

\par\begin{minipage}{\linewidth}\raggedright
% item: EX-I-07-15
% key: KEY-EX-I-07-15
% mdparagraph:P190:begin
\textbf{I.7.15. Let the men swim in the sea.}
% mdparagraph:P190:end

% mdparagraph:P191:begin
\textbf{Model:} \textbf{Swimmen ðā męnn on ðǣm sǣ.}
% mdparagraph:P191:end

% mdparagraph:P192:begin
\textbf{Target:} third plural present subjunctive and the dative location phrase with selected masculine \textbf{sǣ}. \textbf{Męnn} is nominative plural here, although other cells share the spelling. \textbf{Swimmað} could be an assertion or a command in other settings, but would not be the selected third-person wish.
% mdparagraph:P192:end

\end{minipage}\par\addvspace{4pt}

\par\begin{minipage}{\linewidth}\raggedright
% item: EX-I-07-16
% key: KEY-EX-I-07-16
% mdparagraph:P193:begin
\textbf{I.7.16. Let the tree bear fruit.}
% mdparagraph:P193:end

% mdparagraph:P194:begin
\textbf{Model:} \textbf{Bere ðæt trēo wæstm.}
% mdparagraph:P194:end

% mdparagraph:P195:begin
\textbf{Target:} singular present subjunctive \textbf{bere}, nominative singular tree, and accusative fruit in the produce-bearing sense. \textbf{Wæstmas} is defensible for individual fruits. \textbf{Birð} would assert that the tree bears fruit; \textbf{berað} would also lose the requested singular subjunctive. Do not make the subject plural from the collective interpretation of fruit.
% mdparagraph:P195:end

\end{minipage}\par\addvspace{4pt}

\par\begin{minipage}{\linewidth}\raggedright
% item: EX-I-07-17
% key: KEY-EX-I-07-17
% mdparagraph:P196:begin
\textbf{I.7.17. He says likewise, Let the light be for a sign.}
% mdparagraph:P196:end

% mdparagraph:P197:begin
\textbf{Model:} \textbf{Hē sægð swilce: “Bēo ðæt lēoht tō tācne.”}
% mdparagraph:P197:end

% mdparagraph:P198:begin
\textbf{Target:} retain the statement and the quoted wish. \textbf{Swilce} is the adverb “likewise”; \textbf{ēac swilce} can also convey “also, likewise.” \textbf{Bēo} is third singular present subjunctive with subject \textbf{ðæt lēoht}; \textbf{sīe} is an alternative. \textbf{Tō tācne} expresses the intended function, with dative \textbf{tācne}. The quotation is direct speech, not an unintroduced indirect content clause. A second-person command would misidentify the light’s role.
% mdparagraph:P198:end

\end{minipage}\par\addvspace{4pt}

\reserve{7\baselineskip}
% mdparagraph:P199:begin
\section*{D. Binding the animals}
% mdparagraph:P199:end

% mdparagraph:P200:begin
This complete account is new teaching material, not an adapted source passage. The wording is reproduced here so that the interpretation can be checked against the Old English.
% mdparagraph:P200:end

% mdparagraph:P201:begin
\par\begin{minipage}{\linewidth}\raggedright
\readingtext{Twēgen męnn wǣron on ðǣm strande. Hīe hæfdon fīf dēor. Se ǣresta mann cwæð: “Gif ðū mē helpe, wē magon ðā dēor bindan.” Se ōðer mann cwæð: “Ic wille ðē helpan. Binden wē ðā dēor.” Hīe bundon ðā dēor, and se ǣresta mann sealde his gefēran męte. Ðā dēor wǣron bunden, and ðā męnn hæfdon męte.}
\end{minipage}\par
% mdparagraph:P201:end

\par\begin{minipage}{\linewidth}\raggedright
% item: NEW-U07-D01
% key: KEY-NEW-U07-D01
% mdparagraph:P202:begin
\textbf{D01. Complete translation.}
% mdparagraph:P202:end

% mdparagraph:P203:begin
Two men were on the shore. They had five animals. The first man said, “If you help me, we can bind the animals.” The second man said, “I am willing to help you. Let us bind the animals.” They bound the animals, and the first man gave his companion food. The animals were bound, and the men had food.
% mdparagraph:P203:end

% mdparagraph:P204:begin
\textbf{Target:} preserve the continuous account: two men, five animals, the request, the agreement to help, the shared action, the gift, and the final result. “I want to help you” or “I will help you” can translate \textbf{Ic wille ðē helpan}, provided willingness or undertaking is understood. \textbf{Helpe} makes the first condition tentative; ordinary English “if you help me” can convey that without adding “should.” The account does not say that the animals are killed, that they are the food, or that the first man binds them alone. Do not add those events to make a different story.
% mdparagraph:P204:end

\end{minipage}\par\addvspace{4pt}

\par\begin{minipage}{\linewidth}\raggedright
% item: NEW-U07-D02
% key: KEY-NEW-U07-D02
% mdparagraph:P205:begin
\textbf{D02. What the two speeches do.}
% mdparagraph:P205:end

% mdparagraph:P206:begin
In the first speech, \textbf{Gif ðū mē helpe} is the condition: the first man tentatively seeks the second man’s help. \textbf{Helpe} is present subjunctive, second singular. \textbf{Wē magon ðā dēor bindan} claims their ability if that help is given; \textbf{magon} is present indicative, first plural, with infinitive \textbf{bindan}.
% mdparagraph:P206:end

% mdparagraph:P207:begin
In the second speech, \textbf{Ic wille ðē helpan} expresses the second man’s willingness to help the first. \textbf{Wille} is first singular present indicative in this assertion of volition. \textbf{Binden wē ðā dēor} then urges a shared action, “Let us bind the animals”: \textbf{binden} is present subjunctive, first plural.
% mdparagraph:P207:end

% mdparagraph:P208:begin
\textbf{Target:} distinguish contingent help, ability, willingness, and exhortation. The identical spelling of indicative and subjunctive \textbf{wille} does not erase the explicit speaker’s undertaking. Neither \textbf{magon} nor \textbf{wille} is a passive auxiliary, and \textbf{binden} here is finite rather than the infinitive \textbf{bindan}.
% mdparagraph:P208:end

\end{minipage}\par\addvspace{4pt}

\par\begin{minipage}{\linewidth}\raggedright
% item: NEW-U07-D03
% key: KEY-NEW-U07-D03
% mdparagraph:P209:begin
\textbf{D03. Datives and the companion.}
% mdparagraph:P209:end

% mdtable:T15:begin
\par\begingroup
\fontsize{10.5}{13.2}\selectfont
\setlength{\LTleft}{0pt}\setlength{\LTright}{\fill}
\noindent
\begin{tabular}{@{}>{\raggedright\arraybackslash}p{28.61mm}>{\raggedright\arraybackslash}p{58.87mm}>{\raggedright\arraybackslash}p{30.89mm}@{}}
\toprule
\textbf{Expression} & \textbf{Case and relation} & \textbf{Referent in this account} \\
\midrule
mē in ðū mē helpe & dative singular, governed by helpan & the first man, speaking \\
ðē in Ic wille ðē helpan & dative singular, governed by helpan & the first man, now addressed by the second \\
his gefēran in sealde his gefēran męte & gefēran is dative singular, the gift-recipient; his is a dependent genitive possessor & the first man’s companion, the second man \\
\bottomrule
\end{tabular}
\par\endgroup

% mdtable:T15:end
% mdparagraph:P210:begin
\textbf{Target:} two kinds of dative relation: the person helped with \textbf{helpan}, and the person receiving the gift with \textbf{sęllan}. \textbf{Męte} is the accusative gift, not another recipient. The immediately stated giver \textbf{se ǣresta mann} and the speaker sequence make \textbf{his} naturally refer to the first man. \textbf{His} is not inherently a reflexive form that must always refer to a clause’s subject; the identification comes from this narrative.
% mdparagraph:P210:end

\end{minipage}\par\addvspace{4pt}

\par\begin{minipage}{\linewidth}\raggedright
% item: NEW-U07-D04
% key: KEY-NEW-U07-D04
% mdparagraph:P211:begin
\textbf{D04. The final passive and its cause.}
% mdparagraph:P211:end

% mdparagraph:P212:begin
The passive is \textbf{Ðā dēor wǣron bunden}, “The animals were bound.” \textbf{Ðā dēor} is nominative plural; \textbf{wǣron} is third plural preterite indicative; \textbf{bunden} is the uninflected past participle in the verbal predicate.
% mdparagraph:P212:end

% mdparagraph:P213:begin
The earlier \textbf{Hīe bundon ðā dēor}, “They bound the animals,” establishes the result. \textbf{Hīe} refers to the two men who have just agreed to bind them together. The final passive omits an agent phrase, but the preceding active clause supplies the agents and the event. \textbf{Wǣron} suits the resulting state here; that contextual reading is not an absolute rule that every \textbf{wesan} passive must describe only a state.
% mdparagraph:P213:end

\end{minipage}\par\addvspace{4pt}

\reserve{7\baselineskip}
% mdparagraph:P214:begin
\section*{R. The fifth day of Creation}
% mdparagraph:P214:end

% mdparagraph:P215:begin
The following is the complete selected paragraph of Cook’s \emph{First Book in Old English}, second edition, p. 126, one-based PDF page 142. Its boundaries are \textbf{God cwæð ēac swilce} through \textbf{se fīfta dæg}, editorially counted prose lines 5–14 on that page, checked against its printed line anchors. These are not a separately printed consecutive line-number series. [7]
% mdparagraph:P215:end

% mdparagraph:P216:begin
The source’s words, spellings, quantities, endings, and punctuation are retained. Line wrapping and the purely line-end hyphens in \textbf{swim-/\allowbreak{}mendu} and \textbf{styriend-/\allowbreak{}licu} are joined. Quotation marks are represented by modern curved double quotation marks. The original superscript note cues are detached from the running text; all their numbers, attachments, and wording are retained below. In particular, \textbf{męrgen} and \textbf{styriendlicu} have not been silently regularised.
% mdparagraph:P216:end

% mdparagraph:P217:begin
\par\begin{minipage}{\linewidth}\raggedright
\readingtext{God cwæð ēac swilce, “Tēon nū ðā wæteru forð swimmendu cynn cucu on līfe, and flēogendu cynn ofer eorðan under ðǣre heofonan fæstnesse.” And God gescēop ðā ðā miclan hwalas, and eall libbendu fisccynn and styriendlicu, ðe ðā wæteru tugon forð on hiera hīwum, and eall flēogendu cynn æfter hiera cynne; God geseah ðā ðæt hit gōd wæs. And blētsode hīe, ðus cweðende, “Weaxað, and bēoð gemanigfielde, and gefyllað ðǣre sǣ wæteru, and ðā fuglas bēon gemanigfielde ofer eorðan.” And ðā wæs geworden ǣfen and męrgen se fīfta dæg.}
\end{minipage}\par
% mdparagraph:P217:end

\reserve{7\baselineskip}
% mdparagraph:P218:begin
\subsection*{Cook’s original notes}
% mdparagraph:P218:end

% mdtable:T16:begin
\par\begingroup
\fontsize{10.5}{13.2}\selectfont
\setlength{\LTleft}{0pt}\setlength{\LTright}{\fill}
\begin{longtable}{@{}>{\raggedright\arraybackslash}p{26.09mm}>{\raggedright\arraybackslash}p{43.4mm}>{\raggedright\arraybackslash}p{48.88mm}@{}}
\toprule
\textbf{Original note} & \textbf{Attachment} & \textbf{Original wording} \\
\midrule
\endfirsthead
\toprule
\textbf{Original note} & \textbf{Attachment} & \textbf{Original wording} \\
\midrule
\endhead
\midrule
\multicolumn{3}{r}{\footnotesize Continued on the following page.}\\
\endfoot
\bottomrule
\endlastfoot
4 & ēac swilce & Ēac swilce = etiam. \\
5 & first forð & Producant = tēon . . . forð. \\
6 & cucu & See cucu. \\
7 & līfe & See līf. \\
8 & first flēogendu & See flēogan, and 61. \\
9 & first ðā after gescēop & Adverb; see 84. 1. \\
10 & hwalas & See hwæl. \\
11 & styriendlicu & Lat. motabilem. \\
12 & ðe & Acc. \\
13 & ðā before wæteru tugon & Nom. plur. \\
14 & tugon & See tēon. \\
15 & forð after tugon & Tugon forð = producerunt. \\
16 & blētsode & See blētsian, and 33. \\
17 & cweðende & See cweðan. \\
18 & Weaxað & See weaxan, and 24. \\
19 & first gemanigfielde & Past part. in nom. plur. \\
20 & gefyllað & See gefyllan. \\
21 & bēon & See 193. a. \\
\end{longtable}
\par\endgroup

% mdtable:T16:end
\reserve{7\baselineskip}
% mdparagraph:P219:begin
\subsection*{Local reading help}
% mdparagraph:P219:end

% mdparagraph:P220:begin
This help is part of the reading task. Use it when judging the translation and analyses; it is not a list of constructions to reproduce independently.
% mdparagraph:P220:end

% mdparagraph:P221:begin
\itempara{1}{\textbf{Tēon … forð /\allowbreak{} tugon forð:} “bring forth /\allowbreak{} brought forth.” The first is a present third-plural subjunctive with the waters as subject; the second is a preterite third-plural indicative. The phrase’s bring-forth sense is supplied separately from the earlier draw/\allowbreak{}pull sense of \textbf{tēon}.}
% mdparagraph:P221:end

% mdparagraph:P222:begin
\itempara{2}{\textbf{Swimmendu cynn cucu on līfe /\allowbreak{} flēogendu cynn:} “swimming kinds, alive/\allowbreak{}living /\allowbreak{} flying kinds.” \textbf{Cynn} is neuter accusative plural, the things to be brought forth. The exact participial forms in \textbf{-u} are recognised in this passage; irregular \textbf{cucu} and the alive-expression are supplied. Do not infer a mechanical \textbf{-ende → -endu} production rule for every participle.}
% mdparagraph:P222:end

% mdparagraph:P223:begin
\itempara{3}{\textbf{Ofer eorðan /\allowbreak{} under ðǣre heofonan fæstnesse:} “over the earth /\allowbreak{} under the firmament of heaven.” \textbf{Heofonan} is the genitive singular of the source’s weak feminine \textbf{heofone}; \textbf{fæstnesse} is dative after \textbf{under}. \textbf{Ðǣre} can agree with the genitive \textbf{heofonan} or the dative \textbf{fæstnesse}: its spelling permits both. \textbf{Eorðan} likewise does not distinguish accusative from dative after spatial \textbf{ofer}. Do not demand a uniquely identifiable case where the text does not provide one.}
% mdparagraph:P223:end

% mdparagraph:P224:begin
\itempara{4}{\textbf{Ðā ðā miclan hwalas /\allowbreak{} eall libbendu fisccynn and styriendlicu:} “then the great whales /\allowbreak{} all living and moving kinds of fish.” The first \textbf{ðā} means “then”; the second is the article. \textbf{Fisccynn} is neuter accusative plural. \textbf{Libbendu} and \textbf{styriendlicu} describe it; the latter’s source \textbf{-u} is supplied. \textbf{And styriendlicu} does not require an invented second noun meaning “moving creatures.”}
% mdparagraph:P224:end

% mdparagraph:P225:begin
\itempara{5}{\textbf{Ðe ðā wæteru tugon forð on hiera hīwum:} “which the waters brought forth according to their forms/\allowbreak{}kinds.” Relative \textbf{ðe} is the object; \textbf{ðā wæteru} is the subject, as Cook’s notes 12–13 confirm. The whole relative construction and \textbf{on hiera hīwum} are supplied. \textbf{Hīwum} is dative plural of neuter \textbf{hīw}.}
% mdparagraph:P225:end

% mdparagraph:P226:begin
\itempara{6}{\textbf{Æfter hiera cynne:} “according to their kind.” \textbf{Cynne} is dative; this classificatory use of \textbf{æfter} is supplied.}
% mdparagraph:P226:end

% mdparagraph:P227:begin
\itempara{7}{\textbf{God geseah ðā ðæt hit gōd wæs:} “God then saw that it was good.” The content clause, its word order, and the situation-reference of \textbf{hit} are supplied. Recognising the individual indicative forms does not establish every use of content clauses.}
% mdparagraph:P227:end

% mdparagraph:P228:begin
\itempara{8}{\textbf{And blētsode hīe, ðus cweðende:} “And [God] blessed them, saying thus.” The subject God continues from the preceding clause. Plain \textbf{blētsode} is supplied third singular preterite of \textbf{blētsian}, distinct from the course’s productive \textbf{geblētsian}. \textbf{Cweðende} describes the same subject’s accompanying speech.}
% mdparagraph:P228:end

% mdparagraph:P229:begin
\itempara{9}{\textbf{Weaxað … bēoð … gefyllað /\allowbreak{} ðā fuglas bēon gemanigfielde:} the first verbs address a plural group; the last clause expresses a third-person wish about the birds. \textbf{Gefyllað} means “fill” here, a supplied sense beyond finish/\allowbreak{}accomplish. \textbf{Ðǣre sǣ wæteru} means “the waters of the sea”: \textbf{sǣ} is source feminine genitive singular, with agreeing \textbf{ðǣre}. This is separate from the earlier productive masculine paradigm.}
% mdparagraph:P229:end

% mdparagraph:P230:begin
\itempara{10}{\textbf{And ðā wæs geworden ǣfen and męrgen se fīfta dæg:} “And then evening and morning had come to be, the fifth day.” The intransitive result formula is supplied. It is not an ordinary transitive passive, and it does not by itself establish a general agreement rule or a complete perfect-tense system. Preserve the source spelling \textbf{męrgen}.}
% mdparagraph:P230:end

\par\begin{minipage}{\linewidth}\raggedright
% item: NEW-U07-R01
% key: KEY-NEW-U07-R01
% mdparagraph:P231:begin
\textbf{R01. Complete translation.}
% mdparagraph:P231:end

% mdparagraph:P232:begin
God also said, “Let the waters now bring forth swimming kinds, alive with life, and flying kinds over the earth under the firmament of heaven.” And God then created the great whales, and all living and moving kinds of fish which the waters brought forth according to their forms, and all flying kinds according to their kind; God then saw that it was good. And he blessed them, saying thus, “Increase, and be multiplied, and fill the waters of the sea, and let the birds be multiplied over the earth.” And then evening and morning had come to be, the fifth day.
% mdparagraph:P232:end

% mdparagraph:P233:begin
\textbf{Target:} all stages of the paragraph: the initial wish, creation of the creatures, God’s judgment and blessing, the commands and bird-wish, and the fifth-day close. “Living swimming kinds” can render \textbf{swimmendu cynn cucu on līfe} more smoothly than the deliberately close “alive with life.” “Kinds of flying creatures” can render \textbf{flēogendu cynn}. “According to their kinds” is defensible for the classificatory phrases if the relation is preserved. “Evening and morning came to be” can convey the final formula without importing an English pluperfect system into Old English.
% mdparagraph:P233:end

% mdparagraph:P234:begin
Keep \textbf{God} as the subject of creating, seeing, and blessing. The waters are the subject of \textbf{tugon}; the relative \textbf{ðe} is what they brought forth. Do not turn the birds’ final third-person wish into a second-person command merely because it follows commands. A good translation need not resolve the genuine case ambiguities in the spatial phrases beyond the meaning supplied above.
% mdparagraph:P234:end

\end{minipage}\par\addvspace{4pt}

\par\begin{minipage}{\linewidth}\raggedright
% item: NEW-U07-R02
% key: KEY-NEW-U07-R02
% mdparagraph:P235:begin
\textbf{R02. Four finite forms in their clauses.}
% mdparagraph:P235:end

% mdtable:T17:begin
\par\begingroup
\fontsize{10.5}{13.2}\selectfont
\setlength{\LTleft}{0pt}\setlength{\LTright}{\fill}
\noindent
\begin{tabular}{@{}>{\raggedright\arraybackslash}p{15.78mm}>{\raggedright\arraybackslash}p{47.35mm}>{\raggedright\arraybackslash}p{55.24mm}@{}}
\toprule
\textbf{Source form} & \textbf{Analysis in this passage} & \textbf{Subject/\allowbreak{}addressee and force} \\
\midrule
Tēon & present subjunctive, third plural, of tēon & ðā wæteru; “Let the waters bring forth” \\
Weaxað & second plural imperative of weaxan & the created beings addressed as “you”; “Increase” \\
bēoð & second plural imperative of bēon, passive auxiliary & the same plural addressees; “be multiplied” \\
bēon & present subjunctive, third plural, passive auxiliary & explicit subject ðā fuglas; “let the birds be multiplied” \\
\bottomrule
\end{tabular}
\par\endgroup

% mdtable:T17:end
% mdparagraph:P236:begin
\textbf{Target:} identify finite forms from their clause relations. Initial \textbf{Tēon} is not an infinitive despite the shared spelling; it has the nominative subject \textbf{ðā wæteru} in an independent wish. \textbf{Weaxað} and \textbf{bēoð} share spellings with indicative plural forms, but the blessing’s direct address makes them commands here. Final \textbf{bēon} has third-person birds as its subject, rather than an unexpressed second-person “you.”
% mdparagraph:P236:end

\end{minipage}\par\addvspace{4pt}

\par\begin{minipage}{\linewidth}\raggedright
% item: NEW-U07-R03
% key: KEY-NEW-U07-R03
% mdparagraph:P237:begin
\textbf{R03. The recognised participial forms.}
% mdparagraph:P237:end

% mdtable:T18:begin
\par\begingroup
\fontsize{10.5}{13.2}\selectfont
\setlength{\LTleft}{0pt}\setlength{\LTright}{\fill}
\noindent
\begin{tabular}{@{}>{\raggedright\arraybackslash}p{39.66mm}>{\raggedright\arraybackslash}p{40.61mm}>{\raggedright\arraybackslash}p{38.11mm}@{}}
\toprule
\textbf{Occurrence} & \textbf{Noun modified} & \textbf{Required agreement in this occurrence} \\
\midrule
swimmendu in swimmendu cynn & cynn, object of Tēon … forð & accusative neuter plural \\
first flēogendu in flēogendu cynn & cynn, coordinated object of Tēon … forð & accusative neuter plural \\
libbendu in libbendu fisccynn & fisccynn, object of gescēop & accusative neuter plural \\
second flēogendu in eall flēogendu cynn & cynn, coordinated object of gescēop & accusative neuter plural \\
\bottomrule
\end{tabular}
\par\endgroup

% mdtable:T18:end
% mdparagraph:P238:begin
\textbf{Target:} all four occurrences of the three given spellings. Their nouns are plural although \textbf{cynn} and \textbf{fisccynn} have zero plural endings. “The waters” is not the noun modified by \textbf{swimmendu}, and the nominative subjects do not supply these participles’ cases. Recognise these exact source forms with the supplied clause help; this task does not ask for unaided formation of other participles or other cells of these lemmas.
% mdparagraph:P238:end

\end{minipage}\par\addvspace{4pt}

\par\begin{minipage}{\linewidth}\raggedright
% item: NEW-U07-R04
% key: KEY-NEW-U07-R04
% mdparagraph:P239:begin
\textbf{R04. Gemanigfielde in the blessing.}
% mdparagraph:P239:end

% mdparagraph:P240:begin
Both occurrences are past participles of \textbf{gemanigfieldan}, in verbal passive predicates. In \textbf{bēoð gemanigfielde}, the subject is the unexpressed plural “you” addressed by the command. In \textbf{ðā fuglas bēon gemanigfielde}, the subject is explicit third-person plural \textbf{ðā fuglas}.
% mdparagraph:P240:end

% mdparagraph:P241:begin
\textbf{Target:} distinguish the subjects and the finite auxiliaries while explaining the common plural participial ending. Cook identifies the first \textbf{gemanigfielde} as nominative plural. Its \textbf{-e} is compatible with the permitted plural predicate ending; it does not turn the participle into a finite present verb or an infinitive. An uninflected participle is also permitted in the verbal-predicate system, so the source’s ending is not evidence that every passive must carry full attributive gender agreement. The source itself is retained exactly.
% mdparagraph:P241:end

\end{minipage}\par\addvspace{4pt}

\par\begin{minipage}{\linewidth}\raggedright
% item: NEW-U07-R05
% key: KEY-NEW-U07-R05
% mdparagraph:P242:begin
\textbf{R05. Ðǣre sǣ wæteru and on ðǣm sǣ.}
% mdparagraph:P242:end

% mdparagraph:P243:begin
\textbf{Ðǣre sǣ wæteru} means \textbf{“the waters of the sea.”} \textbf{Wæteru} is the head, accusative neuter plural as the object of \textbf{gefyllað}. \textbf{Ðǣre sǣ} is the dependent genitive singular; this source use of \textbf{sǣ} is feminine, and \textbf{ðǣre} agrees with it.
% mdparagraph:P243:end

% mdparagraph:P244:begin
\textbf{On ðǣm sǣ} means \textbf{“in the sea.”} Here \textbf{on} governs a dative location phrase, and \textbf{ðǣm} agrees with the course’s previously selected masculine \textbf{sǣ}.
% mdparagraph:P244:end

% mdparagraph:P245:begin
\textbf{Target:} nominal dependence versus prepositional location. Do not call \textbf{ðǣre sǣ} dative merely because the same article spelling can be dative elsewhere. Here its relation to \textbf{wæteru} is genitive. Conversely, do not silently substitute the masculine genitive \textbf{sǣs} for the source’s supplied feminine form.
% mdparagraph:P245:end

\end{minipage}\par\addvspace{4pt}

\par\begin{minipage}{\linewidth}\raggedright
% item: NEW-U07-R06
% key: KEY-NEW-U07-R06
% mdparagraph:P246:begin
\textbf{R06. The relative clause and the sequence.}
% mdparagraph:P246:end

% mdparagraph:P247:begin
The relative clause begins at \textbf{ðe} and ends at \textbf{hīwum}:
% mdparagraph:P247:end

% mdparagraph:P248:begin
\par\begin{minipage}{\linewidth}\raggedright
\readingtext{\textbf{ðe ðā wæteru tugon forð on hiera hīwum} — “which the waters brought forth according to their forms.”}
\end{minipage}\par
% mdparagraph:P248:end

% mdparagraph:P249:begin
\textbf{Ðe} is the accusative relative object; \textbf{ðā wæteru} is its clause’s nominative subject. The preceding \textbf{eall libbendu fisccynn and styriendlicu}, “all living and moving kinds of fish,” supplies what is brought forth in this relative clause. The following \textbf{and eall flēogendu cynn æfter hiera cynne} resumes the objects of God’s creating; it need not be pulled inside the relative clause.
% mdparagraph:P249:end

% mdparagraph:P250:begin
God then blessed \textbf{hīe}, “them,” the created beings just described, including the water creatures and birds. The paragraph ends with \textbf{se fīfta dæg}, “the fifth day.”
% mdparagraph:P250:end

% mdparagraph:P251:begin
\textbf{Target:} follow the connected reference with the printed help. \textbf{Hīe} is the object of \textbf{blētsode}, not a new plural subject “they blessed.” The source’s final \textbf{wæs geworden} describes an intransitive coming-to-be; treating it as an ordinary “was made [by someone]” passive would blur the supplied closing formula.
% mdparagraph:P251:end

\end{minipage}\par\addvspace{4pt}

\reserve{7\baselineskip}
% mdparagraph:P252:begin
\section*{X. Cumulative practice}
% mdparagraph:P252:end

\par\begin{minipage}{\linewidth}\raggedright
% item: NEW-U07-X01
% key: KEY-NEW-U07-X01
% mdparagraph:P253:begin
\textbf{X01. The object of gīeman.}
% mdparagraph:P253:end

% mdparagraph:P254:begin
\textbf{Model:} \textbf{Se mann gīemð ðāra fugla.} — “The man rules over the birds.”
% mdparagraph:P254:end

% mdparagraph:P255:begin
\textbf{Target:} the taught sense of \textbf{gīeman} governs a genitive. \textbf{Ðāra fugla} is therefore genitive plural. An accusative \textbf{ðā fuglas} would not satisfy this frame, even though many earlier transitive verbs take accusatives. \textbf{Se mann} remains nominative singular and controls third singular \textbf{gīemð}. A valid paraphrase with another verb does not answer the requested correction of this verb’s government.
% mdparagraph:P255:end

\end{minipage}\par\addvspace{4pt}

\par\begin{minipage}{\linewidth}\raggedright
% item: NEW-U07-X02
% key: KEY-NEW-U07-X02
% mdparagraph:P256:begin
\textbf{X02. Each tree /\allowbreak{} all trees.}
% mdparagraph:P256:end

% mdparagraph:P257:begin
\textbf{Models:} \textbf{ǣlc trēo}, “each tree,” nominative neuter singular; \textbf{eall trēowu}, “all trees,” accusative neuter plural.
% mdparagraph:P257:end

% mdparagraph:P258:begin
\textbf{Target:} ordinary “each/\allowbreak{}every” here selects a singular noun phrase; \textbf{eall} agrees with the plural trees. Both modifiers have zero endings in these particular cells, but the noun’s \textbf{trēo : trēowu} distinction and the requested syntactic roles remain. Do not infer a new plural meaning of \textbf{ǣlc} merely because its complete morphological table has plural cells.
% mdparagraph:P258:end

\end{minipage}\par\addvspace{4pt}

\par\begin{minipage}{\linewidth}\raggedright
% item: NEW-U07-X03
% key: KEY-NEW-U07-X03
% mdparagraph:P259:begin
\textbf{X03. “He can help the man, and he wants to give the birds food.”}
% mdparagraph:P259:end

% mdparagraph:P260:begin
\textbf{Model:} \textbf{Hē mæg helpan ðǣm męnn, and hē wille giefan ðǣm fuglum męte.}
% mdparagraph:P260:end

% mdparagraph:P261:begin
\textbf{Target:} indicative capability and volition with their infinitives. \textbf{Ðǣm męnn} is dative singular, governed by \textbf{helpan}. With \textbf{giefan}, \textbf{ðǣm fuglum} is the dative plural recipient and \textbf{męte} the accusative singular gift. \textbf{Mæg} and \textbf{wille} have singular subjects; the infinitives keep their own argument relations.
% mdparagraph:P261:end

% mdparagraph:P262:begin
An order such as \textbf{Hē mæg ðǣm męnn helpan} is defensible. A suitable \textbf{sęllan} equivalent can preserve the giving meaning, though the requested analysis of \textbf{giefan} is still its dative recipient plus accusative gift. Do not replace \textbf{ðǣm męnn} with an accusative simply because \textbf{helpan} follows a modal, or call \textbf{wille} subjunctive from its spelling alone.
% mdparagraph:P262:end

\end{minipage}\par\addvspace{4pt}

\par\begin{minipage}{\linewidth}\raggedright
% item: NEW-U07-X04
% key: KEY-NEW-U07-X04
% mdparagraph:P263:begin
\textbf{X04. Change the naming subject.}
% mdparagraph:P263:end

% mdparagraph:P264:begin
\textbf{Ic hāte ðā fæstnesse Heofon.} — “I call the firmament Heaven.”
% mdparagraph:P264:end

% mdparagraph:P265:begin
\textbf{Hē hǣtt ðā fæstnesse Heofon.} — “He calls the firmament Heaven.”
% mdparagraph:P265:end

% mdparagraph:P266:begin
\textbf{Target:} the person change requires \textbf{hāte → hǣtt}, first to third singular present indicative. The naming frame stays unchanged: \textbf{ðā fæstnesse} is the accusative entity named and \textbf{Heofon} the accusative assigned name. \textbf{Hǣtt} is active here; do not replace it with the distinct inherited passive \textbf{hātte} or make the name nominative merely because English uses a bare complement.
% mdparagraph:P266:end

\end{minipage}\par\addvspace{4pt}

\par\begin{minipage}{\linewidth}\raggedright
% item: NEW-U07-X05
% key: KEY-NEW-U07-X05
% mdparagraph:P267:begin
\textbf{X05. Command and wish: “be multiplied.”}
% mdparagraph:P267:end

% mdparagraph:P268:begin
\textbf{Second-plural command:} \textbf{Bēoð gemanigfielde.} — “Be multiplied,” addressed to more than one person.
% mdparagraph:P268:end

% mdparagraph:P269:begin
\textbf{Third-plural wish:} \textbf{Bēon hīe gemanigfielde.} — “Let them be multiplied.”
% mdparagraph:P269:end

% mdparagraph:P270:begin
\textbf{Target:} \textbf{bēoð → bēon} changes the finite auxiliary from second plural imperative to third plural present subjunctive. \textbf{Hīe} makes the wish’s third-person subject explicit. It may be omitted if an independently clear context supplies that subject; there is no rule that every subjunctive must have an overt pronoun. \textbf{Gemanigfielde} stays the same because the participle carries neither finite person nor mood. Uninflected \textbf{gemanigfield} is a defensible predicate alternative in either clause, provided the auxiliary distinction remains.
% mdparagraph:P270:end

\end{minipage}\par\addvspace{4pt}

\par\begin{minipage}{\linewidth}\raggedright
% item: NEW-U07-X06
% key: KEY-NEW-U07-X06
% mdparagraph:P271:begin
\textbf{X06. Purpose and accomplishment.}
% mdparagraph:P271:end

% mdparagraph:P272:begin
\textbf{Model:} \textbf{Ðā gefēran cōmon tō helpanne. Ðæt weorc wæs gefręmed.}
% mdparagraph:P272:end

% mdparagraph:P273:begin
“The companions came to help. The work was accomplished.”
% mdparagraph:P273:end

% mdparagraph:P274:begin
\textbf{Target:} the selected inflected purpose infinitive \textbf{tō helpanne}, followed by a preterite passive with \textbf{wesan}. The companions are the subject of \textbf{cōmon} and the understood helpers. \textbf{Ðæt weorc} is nominative neuter singular, \textbf{wæs} third singular preterite indicative, and \textbf{gefręmed} the past participle of \textbf{gefręmman}. A plain infinitive may have uses elsewhere, but \textbf{tō helpan} does not supply the expressly requested inflected form. A \textbf{weorðan} passive can be valid language but misses the chosen auxiliary in this task. The two sentences do not specify an agent phrase for the accomplishment, so none needs to be invented.
% mdparagraph:P274:end

\end{minipage}\par\addvspace{4pt}

\clearpage
% mdparagraph:P275:begin
\section*{Sources}
% mdparagraph:P275:end

\begingroup\fontsize{10}{12.8}\selectfont
% mdparagraph:P276:begin
Page references distinguish the printed page from the one-based page number of the registered PDF. Grammatical explanations and all independent-composition models in this key are new teaching material based on the checked forms and constructions; only the identified Cook prompts, Creation paragraph, and original notes are source transcriptions.
% mdparagraph:P276:end

% mdparagraph:P277:begin
[1] A. S. Cook, \emph{A First Book in Old English}, second edition (1895), §§93–95, printed pp. 54–55 /\allowbreak{} PDF 70–71, for the subjunctive framework; the corresponding lexical systems in §§102–141, printed pp. 59–80 /\allowbreak{} PDF 75–96, for strong, weak, preterite-present, and anomalous forms. Especially §§114, 122, 126, 138–139 distinguish the lexical stems used in A01–A07. The five newly introduced verbs also use their checked vocabulary entries: \textbf{gemanigfieldan} p. 293 /\allowbreak{} PDF 309, \textbf{gīeman} p. 296 /\allowbreak{} PDF 312, \textbf{rīcsian} p. 309 /\allowbreak{} PDF 325, \textbf{tōtwǣman} p. 316 /\allowbreak{} PDF 332, \textbf{wyrcean} p. 324 /\allowbreak{} PDF 340.
% mdparagraph:P277:end

% mdparagraph:P278:begin
[2] Cook, §§55–62, printed pp. 38–41 /\allowbreak{} PDF 54–57, especially §57.3 for participial adjective variants and §§61–62 for participle inflection; §179, printed pp. 99–100 /\allowbreak{} PDF 115–116, for adjectival agreement and predicative use. These support the distinction between an attributive agreement cell and permitted verbal-predicate endings.
% mdparagraph:P278:end

% mdparagraph:P279:begin
[3] Cook, §§187–188, printed pp. 101–102 /\allowbreak{} PDF 117–118, for passive auxiliaries and the limits of transferring modern tense distinctions.
% mdparagraph:P279:end

% mdparagraph:P280:begin
[4] Cook, §§192–198, printed pp. 103–105 /\allowbreak{} PDF 119–121, for assertions, independent wishes, conditions, and commands. The unreal nonpast sentence in M06 is new, expressly scaffolded teaching material, not a quotation from Cook.
% mdparagraph:P280:end

% mdparagraph:P281:begin
[5] Cook, \emph{Exercises in Old English} (1895): I.4.18, printed p. 2 /\allowbreak{} PDF 10; I.6.1–16, printed p. 3 /\allowbreak{} PDF 11; I.7.1–10, printed p. 3 /\allowbreak{} PDF 11; I.7.11 crossing printed pp. 3–4 /\allowbreak{} PDF 11–12; I.7.12–17, printed p. 4 /\allowbreak{} PDF 12. Original English wording, parentheses, and square brackets are preserved; source line wrapping is removed and the apostrophe in \textbf{God's} uses the course’s plain-text transcription convention.
% mdparagraph:P281:end

% mdparagraph:P282:begin
[6] Cook, \emph{First Book}, §156(f), printed p. 91 /\allowbreak{} PDF 107, and vocabulary \textbf{gīeman}, printed p. 296 /\allowbreak{} PDF 312, for the governed genitive; §173, printed p. 97 /\allowbreak{} PDF 113, for the naming double accusative. The authority use of \textbf{ofer} with accusative is directly illustrated in the Creation passage, printed p. 127 /\allowbreak{} PDF 143, including \textbf{ofer ealle gesceafta}. The selected spatial phrases and function phrases are checked against the Creation text and vocabulary. \textbf{Ǣlc} has the printed long initial vowel in the \emph{First Book} vocabulary, p. 274 /\allowbreak{} PDF 290, and in \emph{Exercises}, p. 46 /\allowbreak{} PDF 54.
% mdparagraph:P282:end

% mdparagraph:P283:begin
[7] Cook, \emph{First Book}, Reader I, “The Creation,” fifth-day paragraph, printed p. 126 /\allowbreak{} PDF 142, editorially counted prose lines 5–14 and original notes 4–21. The reading help is new editorial support; it is kept separate from the exact source text and Cook’s original apparatus.
% mdparagraph:P283:end

\endgroup
\end{document}
